% Rozdzial 18: Testowanie, debugowanie i profilowanie w C++20.
% Wylacznie tresc rozdzialu; bez \documentclass i \begin{document}.
% Zgodny ze skryptem: listings, amsmath, amssymb, booktabs, tabularx,
% enumitem, hyperref oraz makra \code, \term, \nota.
% Przyklady terminalowe: zasadniczo Linux/macOS z GCC/Clang;
% roznice na Windows/MSVC zaznaczono osobno.
% Uzycie: % Rozdzial 18: Testowanie, debugowanie i profilowanie w C++20.
% Wylacznie tresc rozdzialu; bez \documentclass i \begin{document}.
% Zgodny ze skryptem: listings, amsmath, amssymb, booktabs, tabularx,
% enumitem, hyperref oraz makra \code, \term, \nota.
% Przyklady terminalowe: zasadniczo Linux/macOS z GCC/Clang;
% roznice na Windows/MSVC zaznaczono osobno.
% Uzycie: % Rozdzial 18: Testowanie, debugowanie i profilowanie w C++20.
% Wylacznie tresc rozdzialu; bez \documentclass i \begin{document}.
% Zgodny ze skryptem: listings, amsmath, amssymb, booktabs, tabularx,
% enumitem, hyperref oraz makra \code, \term, \nota.
% Przyklady terminalowe: zasadniczo Linux/macOS z GCC/Clang;
% roznice na Windows/MSVC zaznaczono osobno.
% Uzycie: % Rozdzial 18: Testowanie, debugowanie i profilowanie w C++20.
% Wylacznie tresc rozdzialu; bez \documentclass i \begin{document}.
% Zgodny ze skryptem: listings, amsmath, amssymb, booktabs, tabularx,
% enumitem, hyperref oraz makra \code, \term, \nota.
% Przyklady terminalowe: zasadniczo Linux/macOS z GCC/Clang;
% roznice na Windows/MSVC zaznaczono osobno.
% Uzycie: \input{rozdzial_18_testowanie_debugowanie_profilowanie_cpp.tex}

\chapter{Testowanie, debugowanie i profilowanie C++}
\label{chap:cpp18-test-debug-profile}

\section{Trzy różne pytania o program}

Po napisaniu i skompilowaniu programu nie wiemy jeszcze, czy jest on
poprawny, bezpieczny i wydajny. Każdą z tych własności sprawdza się
innymi metodami:
\begin{description}[style=nextline]
\item[Testowanie (\emph{testing})]
Czy program daje oczekiwane wyniki na wybranych danych oraz czy
spełnia wymagane kontrakty? Testy wykrywają błędy, lecz ich pozytywny
wynik nie jest ogólnym dowodem poprawności.
\item[Debugowanie (\emph{debugging})]
Gdzie i dlaczego program odbiega od oczekiwanego zachowania?
Debugger pozwala obserwować stos wywołań, zmienne, wątki i kolejne
instrukcje.
\item[Profilowanie (\emph{profiling})]
Gdzie program zużywa czas procesora, pamięć i inne zasoby?
Profilowanie wskazuje rzeczywiste wąskie gardła; samo zgadywanie
na podstawie długości kodu często prowadzi do błędnych optymalizacji.
\end{description}

Istotne jest także rozróżnienie \emph{testu poprawności} od
\emph{benchmarku}: pierwszy sprawdza warunek logiczny, drugi mierzy
wielkość fizyczną (np. czas, przepustowość, alokacje). Szybki, lecz
błędny algorytm nie jest dobrym rozwiązaniem.

\nota{W~tym rozdziale podstawowym standardem jest C++20. GoogleTest,
Google Benchmark, GDB, LLDB, Valgrind i~narzędzia \code{perf} nie są
częścią standardowej biblioteki C++; wymagają odpowiednich instalacji.
Polecenia zależą od systemu operacyjnego i~kompilatora.}

\section{Strategia testowania i testowalny projekt}
\label{sec:cpp18-strategy}

\subsection{Poziomy testów}

\begin{center}
\small
\begin{tabularx}{\textwidth}{@{}p{0.29\textwidth}X@{}}
\toprule
Rodzaj testu & Co sprawdza i kiedy jest przydatny?\\
\midrule
Jednostkowy (\emph{unit}) & Mała, izolowana funkcja lub klasa; szybka informacja zwrotna.\\
Integracyjny & Współdziałanie komponentów, np. parsera i~bazy danych.\\
Systemowy / end-to-end & Cały scenariusz użytkownika, często z~rzeczywistym wejściem/wyjściem.\\
Regresyjny & Przypadek, który wcześniej ujawniał błąd; nie powinien powrócić.\\
Własnościowy / losowy & Niezmienniki dla wielu zestawów danych, nie tylko pojedyncze wyniki.\\
Wydajnościowy & Czas, przepustowość, opóźnienie, pamięć; wymaga kontrolowanego środowiska.\\
\bottomrule
\end{tabularx}
\end{center}

Dobry test powinien być \emph{powtarzalny}, \emph{niezależny} od innych
testów i~możliwie mały. W~praktyce wydzielamy obliczenia od efektów
ubocznych (sieci, zegara, plików, losowości), dzięki czemu zależności
możemy zastąpić kontrolowanymi implementacjami.

\subsection{Kontrakty, przypadki brzegowe i błędy numeryczne}

Rozważmy funkcję obliczającą średnią. Jej kontrakt powinien ustalać
zachowanie dla pustego zakresu; poniżej zgłaszamy wyjątek.

\begin{lstlisting}[caption={Prosta funkcja z jawnym kontraktem dla pustego zakresu (C++20).}]
#include <numeric>
#include <span>
#include <stdexcept>

double average(std::span<const double> xs) {
    if (xs.empty())
        throw std::invalid_argument("empty range");
    return std::accumulate(xs.begin(), xs.end(), 0.0)
         / static_cast<double>(xs.size());
}
\end{lstlisting}

Testujemy nie tylko wartości typowe, lecz także rozmiary $0$ i~$1$,
wartości ujemne, powtórzenia, skrajne zakresy oraz dane naruszające
warunki wejściowe. Dla liczb zmiennoprzecinkowych zwykła równość
\code{a == b} nie zawsze jest właściwa: wynik może zależeć od porządku
dodawania i~optymalizacji. Stosujemy tolerancję dostosowaną do skali
i~wymagań zadania, a~nie arbitralne \code{1e-6} dla każdego algorytmu.

\section{Testy bez zewnętrznego frameworka}

\subsection{\texttt{assert} i \texttt{static\_assert}}

\code{assert(expr)} z~nagłówka \code{<cassert>} sprawdza warunek
w~czasie działania, lecz może zostać usunięte przy zdefiniowanym
\code{NDEBUG}. \code{static_assert} sprawdza warunek w~czasie
kompilacji. W~szczególności zwykłe \code{assert} nie nadaje się do
walidacji niepoprawnych danych wejściowych w~kodzie produkcyjnym.

\begin{lstlisting}[caption={Asercja kompilacyjna i kontrola zachowania funkcji.}]
#include <cassert>
#include <concepts>
#include <type_traits>

constexpr int square(int x) { return x * x; }
static_assert(square(6) == 36);
static_assert(std::integral<int>);
static_assert(!std::is_copy_constructible_v<void>);

int main() {
    assert(square(7) == 49);
}
\end{lstlisting}

\nota{Przy \code{-DNDEBUG} argument \code{assert} może w~ogóle nie być
wykonywany. Nie wolno umieszczać w~nim niezbędnych efektów ubocznych,
np.~\code{assert(++counter == 1)}. \code{static_assert} również nie
zastępuje testów dla danych znanych dopiero podczas wykonania.}

\subsection{Niezależne testy danych wejściowych}

Poniższą funkcję \code{main} umieść pod definicją \code{average}
z~poprzedniego listingu. Test nie wymaga zewnętrznej biblioteki.
Sprawdza wynik, wyjątek i~przypadek jednoelementowy. Kompiluj go
bez \code{-DNDEBUG}, aby wyrażenia \code{assert} były wykonywane.

\begin{lstlisting}[caption={Test funkcji average bez frameworka; dołącz pod jej definicją (C++20).}]
#include <cassert>
#include <cmath>
#include <stdexcept>
#include <vector>

// Zakladamy definicje average z poprzedniego przykladu.
int main() {
    const std::vector<double> xs{1.0, 2.0, 3.0, 4.0};
    assert(std::abs(average(xs) - 2.5) < 1e-12);
    const std::vector<double> one{7.0};
    assert(average(one) == 7.0);

    bool caught = false;
    const std::vector<double> empty;
    try {
        average(empty);
    } catch (const std::invalid_argument&) {
        caught = true;
    }
    assert(caught);
}
\end{lstlisting}

W~realnym projekcie lepiej użyć frameworka, który rejestruje testy,
raportuje niepowodzenia i~pozwala uruchamiać wybrane przypadki.

\section{GoogleTest: testy jednostkowe w praktyce}
\label{sec:cpp18-gtest}

Biblioteka \emph{GoogleTest} (często \code{gtest}) udostępnia makra
\code{TEST}, \code{TEST_F}, \code{EXPECT_*} i~\code{ASSERT_*}.
Assercje \code{EXPECT_*} zapisują niepowodzenie i~pozwalają wykonać
resztę funkcji testowej, podczas gdy \code{ASSERT_*} przerywają
bieżącą funkcję testową. To nie są mechanizmy obsługi wyjątków C++.

\begin{lstlisting}[caption={GoogleTest: wyniki, wyjatki i przyblizona rownosc.}]
// Plik: average_test.cpp
#include <gtest/gtest.h>
#include <stdexcept>
#include <vector>

#include "average.hpp"
// Funkcja average jest zdefiniowana w pliku average.cpp.
TEST(AverageTest, ComputesMean) {
    const std::vector<double> xs{2.0, 4.0, 6.0};
    EXPECT_NEAR(average(xs), 4.0, 1e-12);
}

TEST(AverageTest, RejectsEmptyInput) {
    const std::vector<double> empty;
    EXPECT_THROW(average(empty), std::invalid_argument);
}

TEST(AverageTest, HandlesSingleton) {
    const std::vector<double> xs{7.0};
    ASSERT_EQ(xs.size(), 1u);
    EXPECT_DOUBLE_EQ(average(xs), 7.0);
}
\end{lstlisting}

\code{EXPECT_NEAR(a,b,tol)} sprawdza bezwzględną różnicę, dlatego
niekiedy lepiej użyć własnego testu tolerancji względnej. Makro
\code{EXPECT_DOUBLE_EQ} porównuje liczby z~tolerancją opartą na ULP;
nie należy utożsamiać go z~matematyczną równością liczb rzeczywistych.

\subsection{Fixture i testy parametryzowane}

\emph{Fixture} przygotowuje stan wspólny dla testów. Każdy test
\code{TEST_F} otrzymuje osobny egzemplarz fixture, dzięki czemu zmiany
jednego testu nie powinny wpływać na inny.

\begin{lstlisting}[caption={Fixture z przygotowanym wektorem.}]
#include <gtest/gtest.h>
#include <algorithm>
#include <vector>

class VectorTest : public ::testing::Test {
protected:
    std::vector<int> values{3, 1, 4, 1, 5};
};

TEST_F(VectorTest, SortsAscending) {
    std::sort(values.begin(), values.end());
    EXPECT_TRUE(std::is_sorted(values.begin(), values.end()));
}

TEST_F(VectorTest, StartsFromOriginalState) {
    EXPECT_EQ(values.front(), 3);
    EXPECT_EQ(values.size(), 5u);
}
\end{lstlisting}

Dla wielu par wejście--wynik przydatne są testy parametryzowane:

\begin{lstlisting}[caption={GoogleTest: test parametryzowany z zestawem wartosci.}]
#include <gtest/gtest.h>
#include <tuple>

constexpr int square(int x) { return x * x; }

class SquareTest
    : public ::testing::TestWithParam<std::tuple<int, int>> {};

TEST_P(SquareTest, MatchesExpected) {
    auto [x, expected] = GetParam();
    EXPECT_EQ(square(x), expected);
}

INSTANTIATE_TEST_SUITE_P(
    Values, SquareTest,
    ::testing::Values(
        std::tuple{0, 0},
        std::tuple{2, 4},
        std::tuple{-3, 9}));
\end{lstlisting}

\subsection{Mocki i testowanie współpracy komponentów}

\emph{GoogleMock} pozwala tworzyć kontrolowane implementacje
interfejsów i~sprawdzać wywołania między obiektami. Używaj go tam,
gdzie zachowanie naprawdę zależy od interakcji; nadmierne mockowanie
utrwala przypadkowe szczegóły implementacji.

\begin{lstlisting}[caption={GoogleMock: sprawdzenie pojedynczego wywolania.}]
#include <gmock/gmock.h>
#include <string_view>

struct Sender {
    virtual ~Sender() = default;
    virtual bool send(std::string_view) = 0;
};

struct MockSender : Sender {
    MOCK_METHOD(bool, send, (std::string_view), (override));
};

void notify(Sender& s) { (void)s.send("ready"); }

TEST(NotificationTest, SendsOneMessage) {
    MockSender sender;
    EXPECT_CALL(sender, send("ready"))
        .Times(1).WillOnce(::testing::Return(true));
    notify(sender);
}
\end{lstlisting}

Własne wyjątki, klasy niewspierające kopiowania i~asynchroniczne
wywołania wymagają odpowiedniego projektu fixture/mocków.
Dla operacji wykonywanych w~innym wątku szczególnie istotne jest
kontrolowane oczekiwanie na koniec pracy przed zniszczeniem mocka.

\subsection{CMake i CTest: mały projekt wieloplikowy}
\label{sec:cpp18-cmake-test}

Poniższa organizacja zakłada, że \code{average} ma deklarację w~pliku
\code{average.hpp}, definicję w~\code{average.cpp}, a~testy w~pliku
\code{average_test.cpp}. W~przykładzie \code{average_test.cpp}
należy dodać \verb|#include "average.hpp"| zamiast deklarować
funkcję ręcznie. Plik nagłówkowy może wyglądać następująco:

\begin{lstlisting}[caption={Plik average.hpp: deklaracja API testowanej funkcji.}]
#pragma once
#include <span>
double average(std::span<const double> xs);
\end{lstlisting}

Plik \code{average.cpp} zawiera kod pierwszego przykładu
z~tej części rozdziału (z~dołączonym \code{average.hpp});
w~pliku testów umieszczamy testy z~sekcji~\ref{sec:cpp18-gtest}.

\begin{lstlisting}[language=bash,caption={Struktura projektu z GoogleTest.}]
project/
  CMakeLists.txt
  average.hpp
  average.cpp
  average_test.cpp
\end{lstlisting}

\begin{lstlisting}[language=bash,caption={CMakeLists.txt: biblioteka, testy i automatyczna rejestracja w CTest.}]
cmake_minimum_required(VERSION 3.20)
project(cpp18_tests LANGUAGES CXX)

add_library(core average.cpp)
target_compile_features(core PUBLIC cxx_std_20)
target_include_directories(core PUBLIC ${CMAKE_CURRENT_SOURCE_DIR})

# GoogleTest musi byc zainstalowany lub dostarczony przez menedzer pakietow.
find_package(GTest CONFIG REQUIRED)

add_executable(core_tests average_test.cpp)
target_link_libraries(core_tests PRIVATE core GTest::gtest_main)

enable_testing()
include(GoogleTest)
gtest_discover_tests(core_tests)
\end{lstlisting}

\begin{lstlisting}[language=bash,caption={Budowanie i uruchamianie testow przez CMake/CTest.}]
cmake -S . -B build -DCMAKE_BUILD_TYPE=Debug
cmake --build build -j
ctest --test-dir build --output-on-failure
./build/core_tests --gtest_filter='AverageTest.*'
\end{lstlisting}

\nota{W~systemach z~generatorem wielokonfiguracyjnym (np.~Visual
Studio) wariant budowania wybiera się przez \code{--config Debug}.
Przykład wymaga zainstalowanego GoogleTest; alternatywą jest
\code{FetchContent} z~\emph{przypiętą} wersją zależności, a~nie
pobieranie nieokreślonej wersji podczas każdego budowania.}

\section{Testy własnościowe, generatywne i fuzzing}
\label{sec:cpp18-property}

Sprawdzanie setek wartości jest bardziej użyteczne, jeśli testuje
\emph{własność}, a~nie jedynie listę ręcznie wpisanych wyników.
Na przykład po posortowaniu ciągu musi on być niemalejący i~mieć
tę samą multizbiorowość wartości co przed sortowaniem.

\begin{lstlisting}[caption={Powtarzalny test wlasnosciowy przy uzyciu generatora losowego.}]
#include <algorithm>
#include <cassert>
#include <random>
#include <vector>

int main() {
    std::mt19937 gen(123456u);  // Ustalony seed.
    std::uniform_int_distribution<int> count(0, 300);
    std::uniform_int_distribution<int> value(-1000, 1000);

    for (int trial = 0; trial < 1000; ++trial) {
        std::vector<int> x(count(gen));
        for (int& v : x) v = value(gen);
        auto original = x;
        std::sort(x.begin(), x.end());
        assert(std::is_sorted(x.begin(), x.end()));
        std::sort(original.begin(), original.end());
        assert(x == original);
    }
}
\end{lstlisting}

W~testach losowych przechowuj \emph{seed} i~minimalny przykład,
który ujawnił błąd. Warto również stosować \emph{shrinking}:
automatyczne zmniejszanie wejścia do możliwie prostego kontrprzykładu.

\emph{Fuzzing} automatycznie generuje wejścia, kierując się m.in.
pokryciem kodu. W~Clangu narzędzie \emph{libFuzzer} może wywoływać
poniższą funkcję z~wieloma binarnymi danymi wejściowymi.

\begin{lstlisting}[caption={Minimalny punkt wejsciowy dla libFuzzer (Clang).}]
#include <cstddef>
#include <cstdint>
#include <vector>

extern "C" int LLVMFuzzerTestOneInput(
    const std::uint8_t* data, std::size_t size) {
    // Podmien na testowany parser lub dekoder.
    std::vector<std::uint8_t> bytes(data, data + size);
    if (bytes.size() > 2 && bytes[0] == 0x42) {
        // Przypadek zalezy od danych; nie wolno zakladac ich poprawnosci.
        (void)bytes[1];
    }
    return 0;
}
\end{lstlisting}

\begin{lstlisting}[language=bash,caption={Budowanie harnessu fuzzingu w Clangu.}]
clang++ -std=c++20 -O1 -g \
  -fsanitize=fuzzer,address,undefined fuzz.cpp -o fuzz
./fuzz -max_total_time=30
\end{lstlisting}

Po zgłoszeniu awarii zapisujemy minimalny przypadek jako test
regresyjny. Fuzzer nie dowodzi braku błędów, a~jego skuteczność zależy
od jakości punktów wejściowych i~stanu pokrycia.

\section{Debugowanie: od komunikatu do przyczyny}
\label{sec:cpp18-debug}

\subsection{Informacje debugowe i rodzaje kompilacji}

Użyteczne opcje GCC/Clang to \code{-g} (symbole debugowe),
\code{-O0} (łatwe śledzenie krokowe) oraz \code{-Og} (optymalizacje
starające się zachować dobrą obserwowalność w~debuggerze). Program
kompilowany z~\code{-O2} może nie mieć osobnej zmiennej ani instrukcji
odpowiadającej każdej linii kodu źródłowego.

\begin{lstlisting}[language=bash,caption={Budowanie programu do debugowania.}]
g++ -std=c++20 -g -O0 -Wall -Wextra -Wpedantic \
    app.cpp -o app_debug

gdb ./app_debug
\end{lstlisting}

Nie wszystkie błędy są odtwarzalne przy \code{-O0}; niektóre pojawiają
się dopiero przy \code{-O2} lub \code{-O3}. Do ich analizy warto
pozostawić \code{-g} także w~wersji zoptymalizowanej i~zachować
niestripowane symbole debugowe.

\subsection{GDB: breakpoint, stos i obserwacja zmiennych}

\begin{lstlisting}[caption={Program do cwiczenia debugowania.}]
#include <stdexcept>
#include <vector>

int element_at(const std::vector<int>& v, std::size_t index) {
    return v.at(index);  // Moze zglosic std::out_of_range.
}

int main() {
    const std::vector<int> v{3, 5, 8};
    return element_at(v, 5);
}
\end{lstlisting}

\begin{lstlisting}[language=bash,caption={Typowa sesja w GDB.}]
(gdb) break element_at        # Punkt zatrzymania.
(gdb) run                     # Uruchom program.
(gdb) info args               # Argumenty funkcji.
(gdb) print index             # Wartosc zmiennej.
(gdb) backtrace               # Stos wywolan.
(gdb) next                    # Nastepna linia bez wchodzenia.
(gdb) step                    # Wejdz do wywolywanej funkcji.
(gdb) catch throw             # Zatrzymaj przy zglaszaniu wyjatku.
(gdb) continue                # Kontynuuj.
\end{lstlisting}

\code{break}, \code{watch}, \code{next}, \code{step},
\code{finish}, \code{frame}, \code{print} i~\code{bt} to podstawowe
narzędzia. \emph{Watchpoint} śledzi zmianę pamięci, ale możliwość
sprzętowego śledzenia zależy od architektury i~debuggera.

\subsection{Wątki i deadlock w debuggerze}

W~aplikacji wielowątkowej sama analiza wątku \code{main} nie
wystarczy. W~GDB:

\begin{lstlisting}[language=bash,caption={Diagnostyka wielowatkowosci w GDB.}]
(gdb) info threads
(gdb) thread 3
(gdb) bt
(gdb) thread apply all bt
(gdb) set print pretty on
\end{lstlisting}

Jeśli wątki zatrzymały się w~\code{mutex::lock()} lub
\code{condition_variable::wait()}, trzeba odtworzyć graf
zależności: który wątek posiada dany mutex i~na co czeka.
Program zakleszczony nie musi zawierać wyścigu danych, więc
ThreadSanitizer nie zastąpi analizy zakleszczeń.

\subsection{LLDB: odpowiednik na macOS i innych platformach}

LLDB oferuje analogiczne operacje:

\begin{lstlisting}[language=bash,caption={Podstawowe polecenia LLDB.}]
(lldb) breakpoint set --name element_at
(lldb) run
(lldb) frame variable
(lldb) thread backtrace
(lldb) thread list
(lldb) thread backtrace all
(lldb) next
(lldb) step
(lldb) continue
\end{lstlisting}

W~VS Code można podłączyć GDB/LLDB poprzez odpowiednie rozszerzenie,
a~następnie korzystać z~breakpointów, podglądu zmiennych i~stosów.
Dokładny adapter oraz konfiguracja \code{launch.json} zależą od
systemu, rozszerzenia i~narzędzi zainstalowanych lokalnie.

\section{Sanitizatory i narzędzia wykrywania błędów}
\label{sec:cpp18-sanitizers}

Sanitizatory dodają do programu kod kontrolujący określone klasy
błędów w~czasie wykonywania. W~odróżnieniu od debuggera mogą
wykryć nieprawidłową operację \emph{w~chwili jej wykonania}, zanim
zauważymy błędny rezultat.

\begin{center}
\small
\begin{tabularx}{\textwidth}{@{}p{0.28\textwidth}X@{}}
\toprule
Narzędzie & Główne zastosowanie\\
\midrule
AddressSanitizer (ASan) & Wyjście poza granice, use-after-free, część przypadków wycieków.\\
UndefinedBehaviorSanitizer (UBSan) & Np. przepełnienie \code{signed int}, błędne przesunięcia, null dereference, alignment.\\
ThreadSanitizer (TSan) & Wiele wyścigów danych i~część problemów synchronizacji.\\
MemorySanitizer (MSan) & Odczyty niezainicjalizowanej pamięci; wymaga specjalnych warunków instrumentacji.\\
Valgrind Memcheck & Wykrywanie wybranych błędów pamięci i~wycieków (na wspieranych systemach).\\
\bottomrule
\end{tabularx}
\end{center}

\subsection{AddressSanitizer i UndefinedBehaviorSanitizer}

Poniższy program jest \emph{celowo niepoprawny}; dostęp do elementu
po zwolnieniu pamięci jest niezdefiniowanym zachowaniem. Nie należy
interpretować jego wyniku bez instrumentacji.

\begin{lstlisting}[caption={Celowo bledny program: use-after-free.}]
int main() {
    int* p = new int(42);
    delete p;
    return *p;  // UB: odczyt z nieistniejacego obiektu.
}
\end{lstlisting}

\begin{lstlisting}[language=bash,caption={Instrumentacja ASan i UBSan; GCC/Clang.}]
g++ -std=c++20 -O1 -g -fno-omit-frame-pointer \
    -fsanitize=address,undefined -fno-sanitize-recover=all \
    bad.cpp -o bad_san
./bad_san
\end{lstlisting}

ASan i~UBSan można często połączyć. ASan/TSan/MSan nie są natomiast
zwykle łączone ze sobą w~jednym procesie: używamy osobnych
konfiguracji i~uruchomień. Raporty sanitizatorów zależą od wersji
kompilatora, wsparcia runtime i~platformy.

\subsection{ThreadSanitizer i wyścig danych}

Przykładowy wyścig występuje przy równoczesnym modyfikowaniu
zwykłego licznika. To również \emph{celowo błędny} kod:

\begin{lstlisting}[caption={Celowo bledne wspoldzielenie licznika.}]
#include <thread>

int main() {
    int x = 0;
    std::thread t1([&] { ++x; });
    std::thread t2([&] { ++x; });
    t1.join();
    t2.join();
}
\end{lstlisting}

\begin{lstlisting}[language=bash,caption={Oddzielna konfiguracja z ThreadSanitizer.}]
g++ -std=c++20 -O1 -g -pthread \
    -fsanitize=thread race.cpp -o race_tsan
./race_tsan
\end{lstlisting}

\nota{Brak zgłoszenia w~danym uruchomieniu \emph{nie} dowodzi braku
wyścigów danych. TSan ma ograniczenia platformowe, może zgłaszać
wyniki fałszywie dodatnie dla niestandardowej synchronizacji i~nie
gwarantuje wykrywania deadlocków.}

\subsection{Valgrind i testy wycieków}

Na wspieranych systemach Linux przydatna jest kontrola programów
bez kompilacyjnej instrumentacji:

\begin{lstlisting}[language=bash,caption={Przykladowe badanie pamieci przez Valgrind Memcheck.}]
valgrind --tool=memcheck --leak-check=full \
         --show-leak-kinds=all ./app_debug
\end{lstlisting}

Różne narzędzia mają inne modele wykrywania i~inne koszty
wykonania. Wybór zależy od konkretnej klasy błędów; wyniki dobrze
weryfikować na małym, odtwarzalnym przypadku testowym.

\section{Analiza statyczna, ostrzeżenia i pokrycie kodu}
\label{sec:cpp18-static}

Kompilator potrafi zgłosić wiele usterek jeszcze przed wykonaniem
programu. Warto utrzymywać ostrzeżenia w~CI i~unikać wyciszania ich
bez zrozumienia przyczyny.

\begin{lstlisting}[language=bash,caption={Przykladowe ostrzezenia GCC/Clang.}]
g++ -std=c++20 -Wall -Wextra -Wpedantic -Wconversion \
    -Wshadow -Werror -c core.cpp
\end{lstlisting}

Opcje zależą od kompilatora; \code{-Werror} bywa pożyteczne przy
stabilnej bazie kodu, lecz aktualizacja kompilatora może wprowadzić
nowe diagnostyki. \emph{clang-tidy} wykonuje zestaw analiz
statycznych, np. dotyczących nowoczesnych idiomów C++, błędów
potencjalnych i~stylu.

\begin{lstlisting}[language=bash,caption={Analiza projektu ze zbiorem komend kompilacji.}]
cmake -S . -B build -DCMAKE_EXPORT_COMPILE_COMMANDS=ON
clang-tidy -p build src/core.cpp
\end{lstlisting}

\subsection{Pokrycie testami: co liczby naprawdę oznaczają?}

\emph{Line coverage} mierzy wykonane linie, \emph{branch coverage}
wykonane gałęzie, a~\emph{function coverage} wywołane funkcje.
Wysoki procent nie oznacza sprawdzenia wszystkich kombinacji wejść
ani dowodu poprawności. Na przykład test może wykonać gałąź, nie
sprawdzając rezultatu.

\begin{lstlisting}[language=bash,caption={Pokrycie testow z GCC i gcov (przykladowy projekt).}]
g++ -std=c++20 --coverage -O0 -g core.cpp tests.cpp -o tests
./tests
gcov core.cpp
\end{lstlisting}

Dla Clanga dostępny jest osobny zestaw instrumentacji LLVM:
\begin{lstlisting}[language=bash,caption={Clang: instrumentacja i raport pokrycia.}]
clang++ -std=c++20 -O0 -g -fprofile-instr-generate \
    -fcoverage-mapping core.cpp tests.cpp -o tests
LLVM_PROFILE_FILE=tests.profraw ./tests
llvm-profdata merge -sparse tests.profraw -o tests.profdata
llvm-cov report ./tests -instr-profile=tests.profdata
\end{lstlisting}

W~szczególności należy testować \emph{warunki brzegowe}, a~nie
maksymalizować jedynie procent pokrycia.

\section{Benchmarking: jak mierzyć szybkość bez samooszukiwania}
\label{sec:cpp18-benchmark}

\subsection{Pułapki mikrobenchmarków}

Nowoczesny kompilator usuwa obliczenia o~nieobserwowalnych efektach,
wykonuje \emph{constant folding} i~wektoryzuje pętle. Dlatego mierzenie
pustej pętli lub stałego obliczenia może nie mieć związku z~realną
wydajnością. Wynik zależy ponadto od rozgrzania cache, alokacji,
obciążenia systemu, zegarów CPU, temperatury i~danych testowych.

\begin{lstlisting}[caption={Pomiar algorytmu z rzeczywistym wynikiem (bez biblioteki benchmark).}]
#include <chrono>
#include <iostream>
#include <numeric>
#include <vector>

int main() {
    std::vector<int> data(1'000'000, 1);
    using clock = std::chrono::steady_clock;
    long long checksum = 0;

    for (int repetition = 0; repetition < 10; ++repetition) {
        auto start = clock::now();
        const auto result = std::accumulate(
            data.begin(), data.end(), 0LL);
        auto finish = clock::now();
        checksum += result;  // Wynik jest wykorzystany po pomiarze.
        auto us = std::chrono::duration_cast<
            std::chrono::microseconds>(finish - start).count();
        std::cout << "Trial " << repetition << ": " << us << " us\n";
    }
    std::cout << "checksum=" << checksum << '\n';
}
\end{lstlisting}

To ilustracja, nie wzorcowy mikrobenchmark: optymalizator może
ułatwić sobie obliczenia na regularnych danych; dodatkowo
\code{std::cout} poza odcinkiem pomiarowym zmienia warunki
między próbami. Użyteczne benchmarki wymagają wielu powtórzeń,
przygotowanych danych i~analizy rozkładu czasów.

\subsection{Google Benchmark}

\emph{Google Benchmark} automatyzuje wielokrotne uruchamianie i~raportuje
miary czasowe. W~krytycznym fragmencie wykorzystujemy
\code{benchmark::DoNotOptimize}, aby wynik obliczenia nie został
trywialnie usunięty; nie wyłącza to jednak wszystkich optymalizacji
na badanym wyrażeniu.

\begin{lstlisting}[caption={Samodzielny benchmark sumowania wielu elementow.}]
// Plik: benchmark_sum.cpp
#include <benchmark/benchmark.h>
#include <numeric>
#include <vector>

static void BM_Accumulate(benchmark::State& state) {
    const std::vector<int> data(state.range(0), 1);
    for (auto _ : state) {
        const auto sum = std::accumulate(
            data.begin(), data.end(), 0LL);
        benchmark::DoNotOptimize(sum);
    }
    state.SetItemsProcessed(state.iterations() * state.range(0));
}

BENCHMARK(BM_Accumulate)->Arg(1024)->Arg(1 << 20);
BENCHMARK_MAIN();
\end{lstlisting}

\begin{lstlisting}[language={},caption={CMakeLists.txt dla przykładu Google Benchmark.}]
cmake_minimum_required(VERSION 3.20)
project(BenchmarkSum LANGUAGES CXX)
find_package(benchmark CONFIG REQUIRED)
add_executable(bench benchmark_sum.cpp)
target_compile_features(bench PRIVATE cxx_std_20)
target_link_libraries(bench PRIVATE benchmark::benchmark)
\end{lstlisting}

\begin{lstlisting}[language=bash,caption={Budowanie i powtarzanie pomiaru.}]
cmake -S . -B build -DCMAKE_BUILD_TYPE=Release
cmake --build build
./build/bench --benchmark_repetitions=10
./build/bench --benchmark_out=results.json \
              --benchmark_out_format=json
\end{lstlisting}

W~porównaniach równoległych trzeba ponadto ustalić liczbę wątków,
rozmiary bloków, warm-up, powtarzalność danych oraz ewentualne
\emph{false sharing}. Sam najkrótszy pojedynczy czas nie opisuje
stabilnej wydajności.

\section{Profilowanie CPU i pamięci}
\label{sec:cpp18-profiling}

\subsection{Profilowanie próbkowe: Linux \texttt{perf}}

Zamiast mierzyć wyłącznie całkowity czas, profiler pomaga wykrywać
\emph{hot spots}: funkcje zajmujące znaczącą część wykonania.
Przykładowe polecenia systemu Linux:

\begin{lstlisting}[language=bash,caption={Profilowanie CPU przez perf (Linux).}]
g++ -std=c++20 -O2 -g -fno-omit-frame-pointer \
    app.cpp -o app
perf stat ./app
perf record -g ./app
perf report
\end{lstlisting}

\code{perf stat} podaje liczniki zdarzeń udostępniane przez jądro
i~sprzęt, np. cykle i~instrukcje. \code{perf record} zbiera próbki
umożliwiające oglądanie gorących funkcji i~stosu. Dostępne zdarzenia
i~uprawnienia zależą od komputera i~ustawień systemu.

\subsection{Flame graph}

\emph{Flame graph} to agregat stosów wywołań: szerokość prostokąta
odpowiada udziałowi próbek, a~wysokość zagnieżdżeniu wywołań.
Nie jest to dokładna oś czasu. Przykładowy przepływ z~zewnętrznymi
narzędziami FlameGraph:

\begin{lstlisting}[language=bash,caption={Tworzenie flame graph z probek perf (Linux).}]
perf record -F 99 -g -- ./app
perf script | stackcollapse-perf.pl | flamegraph.pl > flame.svg
\end{lstlisting}

Ostatnie dwa skrypty pochodzą z~osobnego projektu FlameGraph i~muszą
być zainstalowane. W~projektach o~dużej liczbie wątków pomagają też
profile opóźnień, I/O i~blokad, a~nie tylko wykorzystania CPU.

\subsection{Alokacje i zużycie pamięci}

Wydajność pamięci to nie tylko liczba bajtów. Ważne są rozkład
alokacji, fragmentacja, lokalność cache i~czas życia obiektów.
Na wspieranych systemach można użyć m.in.:

\begin{lstlisting}[language=bash,caption={Przykladowe pomiary alokacji i pamieci.}]
valgrind --tool=massif ./app
ms_print massif.out.*
/usr/bin/time -v ./app
\end{lstlisting}

\code{massif} służy do profilowania sterty (zgodnie z~możliwościami
narzędzia); \code{/usr/bin/time -v} raportuje miary procesu, np.
\emph{maximum resident set size}, ale nie mówi, które alokacje były
odpowiedzialne. Instrumentacja pamięci również wpływa na czas pracy.

\section{Asembler i związek kodu źródłowego z wydajnością}
\label{sec:cpp18-asm}

Dwa podobne zapisy C++ mogą dać różny kod maszynowy, szczególnie po
włączeniu \code{-O2} lub \code{-O3}. Przy analizie trzeba jednak
uważać, aby nie wyciągać wniosków na podstawie sztucznych przykładów,
w których kompilator może wyeliminować całą pracę.

\begin{lstlisting}[caption={Niewielka funkcja do analizy wygenerowanego kodu.}]
#include <cstddef>

double dot(const double* a, const double* b, std::size_t n) {
    double sum = 0.0;
    for (std::size_t i = 0; i < n; ++i)
        sum += a[i] * b[i];
    return sum;
}
\end{lstlisting}

\begin{lstlisting}[language=bash,caption={Asembler i raport wektoryzacji w GCC (Linux).}]
g++ -std=c++20 -O3 -S -masm=intel dot.cpp -o dot.s

g++ -std=c++20 -O3 -fopt-info-vec-optimized \
    -c dot.cpp -o dot.o
objdump -drC dot.o
\end{lstlisting}

Dla wektoryzacji ważne są: zależności między iteracjami, wyrównanie,
\emph{aliasing} oraz semantyka arytmetyki zmiennoprzecinkowej.
Brak instrukcji SIMD nie musi oznaczać błędu kompilatora; dla
dokładnie zachowanego porządku sumowania kompilator może nie móc
zmienić kolejności operacji. Flaga \code{-ffast-math} pozwala na
istotnie inne przekształcenia numeryczne i~nie powinna być
włączana automatycznie w~programach wymagających ścisłej semantyki.

\subsection{Profile-guided optimization (PGO)}

Kompilator może wykorzystać profile zebrane na reprezentatywnych
danych, aby lepiej rozłożyć kod i~podjąć decyzje optymalizacyjne.
Schemat dla GCC:

\begin{lstlisting}[language=bash,caption={Dwa przebiegi budowania z PGO (GCC).}]
mkdir -p pgo-data
g++ -std=c++20 -O2 -fprofile-generate=pgo-data \
    app.cpp -o app
./app  # Wykonaj reprezentatywny workload.
g++ -std=c++20 -O2 -fprofile-use=pgo-data \
    app.cpp -o app
\end{lstlisting}

PGO nie zastępuje pomiarów: profil z~innego rozkładu danych może
nawet pogorszyć wyniki. Należy też zachować zgodność źródeł, flag
i~miejsc profilu pomiędzy dwoma etapami budowania.

\section{Testowanie kodu wielowątkowego}
\label{sec:cpp18-concurrency-tests}

W~systemach współbieżnych testy wynikowe muszą badać także
\emph{zakończenie pracy}, własność danych i~synchronizację.
Powtórzenie testu wiele razy zwiększa szansę wykrycia błędu, ale
nigdy nie zastępuje rozumowania o~relacjach \emph{happens-before}
(rozdziały~\ref{chap:concurrency} i~\ref{chap:cpp16-advanced-concurrency}).

\begin{lstlisting}[caption={Poprawny test licznika atomowego z jthread (C++20).}]
#include <atomic>
#include <cassert>
#include <thread>
#include <vector>

int main() {
    std::atomic<int> count{0};
    {
        std::vector<std::jthread> workers;
        for (int i = 0; i < 4; ++i) {
            workers.emplace_back([&] {
                for (int k = 0; k < 10000; ++k)
                    count.fetch_add(1, std::memory_order_relaxed);
            });
        }
    } // Destruktory jthread czekaja na wszystkie watki.
    assert(count.load(std::memory_order_relaxed) == 40000);
}
\end{lstlisting}

Do koordynowania momentu rozpoczęcia testu można wykorzystać
\code{std::latch} lub \code{std::barrier} (C++20). W~testach
\emph{timeout} zapobiega nieskończonemu oczekiwaniu w~infrastrukturze
testowej, ale nie może być uznawany za mechanizm anulowania
wątku. Zwracaj uwagę na wyścigi związane z~destrukcją fixture
oraz na kolejność zamykania wątków i~obiektów synchronizacyjnych.

\section{Praktyczny workflow: od błędu do poprawki}

Prześledźmy zalecaną kolejność działania:
\begin{enumerate}
\item \textbf{Odtwórz błąd}: zanotuj wejście, wersję programu, flagi
      kompilacji, środowisko i~opcjonalnie seed.
\item \textbf{Zminimalizuj przypadek}: mały, powtarzalny przykład jest
      łatwiejszy do analizy niż cały system.
\item \textbf{Włącz diagnostykę}: ostrzeżenia, ASan/UBSan, a~dla
      kodu współbieżnego oddzielne uruchomienie TSan.
\item \textbf{Postaw hipotezę}: zidentyfikuj zmienną lub niezmiennik,
      który naruszono; sprawdź stos i~stan programu w~debuggerze.
\item \textbf{Dodaj test regresyjny}: powinien zawodzić przed poprawką
      i~przechodzić po niej.
\item \textbf{Popraw przyczynę}: nie tylko usuń widoczny symptom.
\item \textbf{Zmierz koszt}: jeżeli poprawka ma znaczenie wydajnościowe,
      porównaj benchmarki i~profile przed oraz po zmianie.
\item \textbf{Uruchom pełen zestaw testów}: szczególnie z~różnymi
      konfiguracjami optymalizacji i~kompilatorami.
\end{enumerate}

\section{Automatyzacja w CI i macierz konfiguracji}

W~projektach utrzymywanych długo testy powinny wykonywać się
automatycznie po każdej istotnej zmianie. Praktyczna macierz obejmuje
zwykle konfigurację debugową z~sanitizatorami oraz konfigurację
optymalizowaną z~testami wynikowymi. TSan uruchamiamy w~osobnej
konfiguracji i~na wspieranym środowisku.

\begin{center}
\small
\begin{tabularx}{\textwidth}{@{}p{0.28\textwidth}X@{}}
\toprule
Konfiguracja & Główne zadanie\\
\midrule
Debug, \code{-O0 -g} & Debugowanie źródłowe i~lokalizacja błędów.\\
Debug/san, \code{-O1 -g}, ASan+UBSan & Dostęp do niepoprawnej pamięci i~wybrane UB.\\
Oddzielny build TSan & Wyścigi danych przy konkurencyjnym wykonaniu.\\
Release, \code{-O2} lub \code{-O3} & Testy produkcyjnej konfiguracji i~benchmarki.\\
Clang i GCC & Różne diagnostyki i~różne decyzje optymalizatora.\\
\bottomrule
\end{tabularx}
\end{center}

\nota{Benchmarki są wrażliwe na obciążenie współdzielonych runnerów
CI. W~zwykłym pipeline CI warto przede wszystkim weryfikować
poprawność i~duże regresje; wiarygodne mikrobenchmarki powinny
wykonywać się na kontrolowanym sprzęcie.}

\section{Najczęstsze błędy i wybór narzędzia}

\begin{center}
\small
\begin{tabularx}{\textwidth}{@{}p{0.40\textwidth}X@{}}
\toprule
Problem & Zalecane podejście\\
\midrule
Funkcja daje złe wyniki & Test jednostkowy, debugger, test regresyjny.\\
Awaria tylko dla niektórych danych & Fuzzing, test własnościowy, ASan/UBSan.\\
Podejrzenie wiszącego wskaźnika & ASan, kontrola czasu życia, debugger.\\
Podejrzenie data race & TSan, analiza synchronizacji i~\emph{happens-before}.\\
Program zawiesza się przy wielu wątkach & Stosy wszystkich wątków, graf blokad, timeout testu.\\
Wysokie zużycie pamięci & Massif, profiler alokacji, pomiar RSS.\\
Program jest wolny & \code{perf}/profiler, później benchmark gorącej funkcji.\\
Spadek szybkości po zmianie & Reprodukowalny benchmark, kontrola warunków, profil porównawczy.\\
Wątpliwości co do wektoryzacji & Raport optymalizacji i~asembler, weryfikacja numeryczna.\\
\bottomrule
\end{tabularx}
\end{center}

\section*{Podsumowanie rozdziału}
\addcontentsline{toc}{section}{Podsumowanie rozdziału}

Testy określają, czy otrzymujemy oczekiwane zachowanie; debugger
pozwala prześledzić przyczyny problemu; sanitizatory wykrywają
określone naruszenia reguł wykonywania; profiler i~benchmarki
pomagają podejmować decyzje dotyczące wydajności. Nie istnieje
jedno narzędzie, które zastępuje wszystkie pozostałe.

Praktyczna sekwencja brzmi: \emph{najpierw poprawność, potem
bezpieczeństwo i~reprodukcyjność, następnie pomiary, wreszcie
optymalizacja}. Każda wykryta usterka powinna prowadzić do testu
regresyjnego, a~każda deklarowana poprawa wydajności --- do
porównywalnego pomiaru. Dla systemów wielowątkowych konieczne jest
ponadto świadome rozumienie synchronizacji, czasu życia obiektów
oraz gwarancji modelu pamięci C++.

\section*{Literatura i dokumentacja}
\addcontentsline{toc}{section}{Literatura i dokumentacja}
\begingroup
\sloppy
\begin{itemize}
\item \url{https://google.github.io/googletest/primer.html} ---
      GoogleTest: asercje, fixture i organizacja testów.
\item \url{https://github.com/google/googletest/blob/main/docs/quickstart-cmake.md}
      --- integracja GoogleTest i CMake.
\item \url{https://cmake.org/cmake/help/latest/module/GoogleTest.html}
      --- \code{gtest_discover_tests} i CTest.
\item \url{https://github.com/google/benchmark/blob/main/docs/user_guide.md}
      --- Google Benchmark i stabilne pomiary.
\item \url{https://clang.llvm.org/docs/AddressSanitizer.html} --- ASan.
\item \url{https://clang.llvm.org/docs/UndefinedBehaviorSanitizer.html}
      --- UBSan.
\item \url{https://clang.llvm.org/docs/ThreadSanitizer.html} --- TSan.
\item \url{https://clang.llvm.org/docs/LibFuzzer.html} --- fuzzing.
\item \url{https://sourceware.org/gdb/current/onlinedocs/gdb.html} --- GDB.
\item \url{https://lldb.llvm.org/use/tutorial.html} --- LLDB.
\item \url{https://valgrind.org/docs/manual/manual.html} --- Memcheck i Massif.
\item \url{https://perf.wiki.kernel.org/} --- profiler \code{perf}.
\item Brendan Gregg, \emph{Systems Performance}, wyd. 2,
      Addison-Wesley, 2020.
\end{itemize}
\endgroup


\chapter{Testowanie, debugowanie i profilowanie C++}
\label{chap:cpp18-test-debug-profile}

\section{Trzy różne pytania o program}

Po napisaniu i skompilowaniu programu nie wiemy jeszcze, czy jest on
poprawny, bezpieczny i wydajny. Każdą z tych własności sprawdza się
innymi metodami:
\begin{description}[style=nextline]
\item[Testowanie (\emph{testing})]
Czy program daje oczekiwane wyniki na wybranych danych oraz czy
spełnia wymagane kontrakty? Testy wykrywają błędy, lecz ich pozytywny
wynik nie jest ogólnym dowodem poprawności.
\item[Debugowanie (\emph{debugging})]
Gdzie i dlaczego program odbiega od oczekiwanego zachowania?
Debugger pozwala obserwować stos wywołań, zmienne, wątki i kolejne
instrukcje.
\item[Profilowanie (\emph{profiling})]
Gdzie program zużywa czas procesora, pamięć i inne zasoby?
Profilowanie wskazuje rzeczywiste wąskie gardła; samo zgadywanie
na podstawie długości kodu często prowadzi do błędnych optymalizacji.
\end{description}

Istotne jest także rozróżnienie \emph{testu poprawności} od
\emph{benchmarku}: pierwszy sprawdza warunek logiczny, drugi mierzy
wielkość fizyczną (np. czas, przepustowość, alokacje). Szybki, lecz
błędny algorytm nie jest dobrym rozwiązaniem.

\nota{W~tym rozdziale podstawowym standardem jest C++20. GoogleTest,
Google Benchmark, GDB, LLDB, Valgrind i~narzędzia \code{perf} nie są
częścią standardowej biblioteki C++; wymagają odpowiednich instalacji.
Polecenia zależą od systemu operacyjnego i~kompilatora.}

\section{Strategia testowania i testowalny projekt}
\label{sec:cpp18-strategy}

\subsection{Poziomy testów}

\begin{center}
\small
\begin{tabularx}{\textwidth}{@{}p{0.29\textwidth}X@{}}
\toprule
Rodzaj testu & Co sprawdza i kiedy jest przydatny?\\
\midrule
Jednostkowy (\emph{unit}) & Mała, izolowana funkcja lub klasa; szybka informacja zwrotna.\\
Integracyjny & Współdziałanie komponentów, np. parsera i~bazy danych.\\
Systemowy / end-to-end & Cały scenariusz użytkownika, często z~rzeczywistym wejściem/wyjściem.\\
Regresyjny & Przypadek, który wcześniej ujawniał błąd; nie powinien powrócić.\\
Własnościowy / losowy & Niezmienniki dla wielu zestawów danych, nie tylko pojedyncze wyniki.\\
Wydajnościowy & Czas, przepustowość, opóźnienie, pamięć; wymaga kontrolowanego środowiska.\\
\bottomrule
\end{tabularx}
\end{center}

Dobry test powinien być \emph{powtarzalny}, \emph{niezależny} od innych
testów i~możliwie mały. W~praktyce wydzielamy obliczenia od efektów
ubocznych (sieci, zegara, plików, losowości), dzięki czemu zależności
możemy zastąpić kontrolowanymi implementacjami.

\subsection{Kontrakty, przypadki brzegowe i błędy numeryczne}

Rozważmy funkcję obliczającą średnią. Jej kontrakt powinien ustalać
zachowanie dla pustego zakresu; poniżej zgłaszamy wyjątek.

\begin{lstlisting}[caption={Prosta funkcja z jawnym kontraktem dla pustego zakresu (C++20).}]
#include <numeric>
#include <span>
#include <stdexcept>

double average(std::span<const double> xs) {
    if (xs.empty())
        throw std::invalid_argument("empty range");
    return std::accumulate(xs.begin(), xs.end(), 0.0)
         / static_cast<double>(xs.size());
}
\end{lstlisting}

Testujemy nie tylko wartości typowe, lecz także rozmiary $0$ i~$1$,
wartości ujemne, powtórzenia, skrajne zakresy oraz dane naruszające
warunki wejściowe. Dla liczb zmiennoprzecinkowych zwykła równość
\code{a == b} nie zawsze jest właściwa: wynik może zależeć od porządku
dodawania i~optymalizacji. Stosujemy tolerancję dostosowaną do skali
i~wymagań zadania, a~nie arbitralne \code{1e-6} dla każdego algorytmu.

\section{Testy bez zewnętrznego frameworka}

\subsection{\texttt{assert} i \texttt{static\_assert}}

\code{assert(expr)} z~nagłówka \code{<cassert>} sprawdza warunek
w~czasie działania, lecz może zostać usunięte przy zdefiniowanym
\code{NDEBUG}. \code{static_assert} sprawdza warunek w~czasie
kompilacji. W~szczególności zwykłe \code{assert} nie nadaje się do
walidacji niepoprawnych danych wejściowych w~kodzie produkcyjnym.

\begin{lstlisting}[caption={Asercja kompilacyjna i kontrola zachowania funkcji.}]
#include <cassert>
#include <concepts>
#include <type_traits>

constexpr int square(int x) { return x * x; }
static_assert(square(6) == 36);
static_assert(std::integral<int>);
static_assert(!std::is_copy_constructible_v<void>);

int main() {
    assert(square(7) == 49);
}
\end{lstlisting}

\nota{Przy \code{-DNDEBUG} argument \code{assert} może w~ogóle nie być
wykonywany. Nie wolno umieszczać w~nim niezbędnych efektów ubocznych,
np.~\code{assert(++counter == 1)}. \code{static_assert} również nie
zastępuje testów dla danych znanych dopiero podczas wykonania.}

\subsection{Niezależne testy danych wejściowych}

Poniższą funkcję \code{main} umieść pod definicją \code{average}
z~poprzedniego listingu. Test nie wymaga zewnętrznej biblioteki.
Sprawdza wynik, wyjątek i~przypadek jednoelementowy. Kompiluj go
bez \code{-DNDEBUG}, aby wyrażenia \code{assert} były wykonywane.

\begin{lstlisting}[caption={Test funkcji average bez frameworka; dołącz pod jej definicją (C++20).}]
#include <cassert>
#include <cmath>
#include <stdexcept>
#include <vector>

// Zakladamy definicje average z poprzedniego przykladu.
int main() {
    const std::vector<double> xs{1.0, 2.0, 3.0, 4.0};
    assert(std::abs(average(xs) - 2.5) < 1e-12);
    const std::vector<double> one{7.0};
    assert(average(one) == 7.0);

    bool caught = false;
    const std::vector<double> empty;
    try {
        average(empty);
    } catch (const std::invalid_argument&) {
        caught = true;
    }
    assert(caught);
}
\end{lstlisting}

W~realnym projekcie lepiej użyć frameworka, który rejestruje testy,
raportuje niepowodzenia i~pozwala uruchamiać wybrane przypadki.

\section{GoogleTest: testy jednostkowe w praktyce}
\label{sec:cpp18-gtest}

Biblioteka \emph{GoogleTest} (często \code{gtest}) udostępnia makra
\code{TEST}, \code{TEST_F}, \code{EXPECT_*} i~\code{ASSERT_*}.
Assercje \code{EXPECT_*} zapisują niepowodzenie i~pozwalają wykonać
resztę funkcji testowej, podczas gdy \code{ASSERT_*} przerywają
bieżącą funkcję testową. To nie są mechanizmy obsługi wyjątków C++.

\begin{lstlisting}[caption={GoogleTest: wyniki, wyjatki i przyblizona rownosc.}]
// Plik: average_test.cpp
#include <gtest/gtest.h>
#include <stdexcept>
#include <vector>

#include "average.hpp"
// Funkcja average jest zdefiniowana w pliku average.cpp.
TEST(AverageTest, ComputesMean) {
    const std::vector<double> xs{2.0, 4.0, 6.0};
    EXPECT_NEAR(average(xs), 4.0, 1e-12);
}

TEST(AverageTest, RejectsEmptyInput) {
    const std::vector<double> empty;
    EXPECT_THROW(average(empty), std::invalid_argument);
}

TEST(AverageTest, HandlesSingleton) {
    const std::vector<double> xs{7.0};
    ASSERT_EQ(xs.size(), 1u);
    EXPECT_DOUBLE_EQ(average(xs), 7.0);
}
\end{lstlisting}

\code{EXPECT_NEAR(a,b,tol)} sprawdza bezwzględną różnicę, dlatego
niekiedy lepiej użyć własnego testu tolerancji względnej. Makro
\code{EXPECT_DOUBLE_EQ} porównuje liczby z~tolerancją opartą na ULP;
nie należy utożsamiać go z~matematyczną równością liczb rzeczywistych.

\subsection{Fixture i testy parametryzowane}

\emph{Fixture} przygotowuje stan wspólny dla testów. Każdy test
\code{TEST_F} otrzymuje osobny egzemplarz fixture, dzięki czemu zmiany
jednego testu nie powinny wpływać na inny.

\begin{lstlisting}[caption={Fixture z przygotowanym wektorem.}]
#include <gtest/gtest.h>
#include <algorithm>
#include <vector>

class VectorTest : public ::testing::Test {
protected:
    std::vector<int> values{3, 1, 4, 1, 5};
};

TEST_F(VectorTest, SortsAscending) {
    std::sort(values.begin(), values.end());
    EXPECT_TRUE(std::is_sorted(values.begin(), values.end()));
}

TEST_F(VectorTest, StartsFromOriginalState) {
    EXPECT_EQ(values.front(), 3);
    EXPECT_EQ(values.size(), 5u);
}
\end{lstlisting}

Dla wielu par wejście--wynik przydatne są testy parametryzowane:

\begin{lstlisting}[caption={GoogleTest: test parametryzowany z zestawem wartosci.}]
#include <gtest/gtest.h>
#include <tuple>

constexpr int square(int x) { return x * x; }

class SquareTest
    : public ::testing::TestWithParam<std::tuple<int, int>> {};

TEST_P(SquareTest, MatchesExpected) {
    auto [x, expected] = GetParam();
    EXPECT_EQ(square(x), expected);
}

INSTANTIATE_TEST_SUITE_P(
    Values, SquareTest,
    ::testing::Values(
        std::tuple{0, 0},
        std::tuple{2, 4},
        std::tuple{-3, 9}));
\end{lstlisting}

\subsection{Mocki i testowanie współpracy komponentów}

\emph{GoogleMock} pozwala tworzyć kontrolowane implementacje
interfejsów i~sprawdzać wywołania między obiektami. Używaj go tam,
gdzie zachowanie naprawdę zależy od interakcji; nadmierne mockowanie
utrwala przypadkowe szczegóły implementacji.

\begin{lstlisting}[caption={GoogleMock: sprawdzenie pojedynczego wywolania.}]
#include <gmock/gmock.h>
#include <string_view>

struct Sender {
    virtual ~Sender() = default;
    virtual bool send(std::string_view) = 0;
};

struct MockSender : Sender {
    MOCK_METHOD(bool, send, (std::string_view), (override));
};

void notify(Sender& s) { (void)s.send("ready"); }

TEST(NotificationTest, SendsOneMessage) {
    MockSender sender;
    EXPECT_CALL(sender, send("ready"))
        .Times(1).WillOnce(::testing::Return(true));
    notify(sender);
}
\end{lstlisting}

Własne wyjątki, klasy niewspierające kopiowania i~asynchroniczne
wywołania wymagają odpowiedniego projektu fixture/mocków.
Dla operacji wykonywanych w~innym wątku szczególnie istotne jest
kontrolowane oczekiwanie na koniec pracy przed zniszczeniem mocka.

\subsection{CMake i CTest: mały projekt wieloplikowy}
\label{sec:cpp18-cmake-test}

Poniższa organizacja zakłada, że \code{average} ma deklarację w~pliku
\code{average.hpp}, definicję w~\code{average.cpp}, a~testy w~pliku
\code{average_test.cpp}. W~przykładzie \code{average_test.cpp}
należy dodać \verb|#include "average.hpp"| zamiast deklarować
funkcję ręcznie. Plik nagłówkowy może wyglądać następująco:

\begin{lstlisting}[caption={Plik average.hpp: deklaracja API testowanej funkcji.}]
#pragma once
#include <span>
double average(std::span<const double> xs);
\end{lstlisting}

Plik \code{average.cpp} zawiera kod pierwszego przykładu
z~tej części rozdziału (z~dołączonym \code{average.hpp});
w~pliku testów umieszczamy testy z~sekcji~\ref{sec:cpp18-gtest}.

\begin{lstlisting}[language=bash,caption={Struktura projektu z GoogleTest.}]
project/
  CMakeLists.txt
  average.hpp
  average.cpp
  average_test.cpp
\end{lstlisting}

\begin{lstlisting}[language=bash,caption={CMakeLists.txt: biblioteka, testy i automatyczna rejestracja w CTest.}]
cmake_minimum_required(VERSION 3.20)
project(cpp18_tests LANGUAGES CXX)

add_library(core average.cpp)
target_compile_features(core PUBLIC cxx_std_20)
target_include_directories(core PUBLIC ${CMAKE_CURRENT_SOURCE_DIR})

# GoogleTest musi byc zainstalowany lub dostarczony przez menedzer pakietow.
find_package(GTest CONFIG REQUIRED)

add_executable(core_tests average_test.cpp)
target_link_libraries(core_tests PRIVATE core GTest::gtest_main)

enable_testing()
include(GoogleTest)
gtest_discover_tests(core_tests)
\end{lstlisting}

\begin{lstlisting}[language=bash,caption={Budowanie i uruchamianie testow przez CMake/CTest.}]
cmake -S . -B build -DCMAKE_BUILD_TYPE=Debug
cmake --build build -j
ctest --test-dir build --output-on-failure
./build/core_tests --gtest_filter='AverageTest.*'
\end{lstlisting}

\nota{W~systemach z~generatorem wielokonfiguracyjnym (np.~Visual
Studio) wariant budowania wybiera się przez \code{--config Debug}.
Przykład wymaga zainstalowanego GoogleTest; alternatywą jest
\code{FetchContent} z~\emph{przypiętą} wersją zależności, a~nie
pobieranie nieokreślonej wersji podczas każdego budowania.}

\section{Testy własnościowe, generatywne i fuzzing}
\label{sec:cpp18-property}

Sprawdzanie setek wartości jest bardziej użyteczne, jeśli testuje
\emph{własność}, a~nie jedynie listę ręcznie wpisanych wyników.
Na przykład po posortowaniu ciągu musi on być niemalejący i~mieć
tę samą multizbiorowość wartości co przed sortowaniem.

\begin{lstlisting}[caption={Powtarzalny test wlasnosciowy przy uzyciu generatora losowego.}]
#include <algorithm>
#include <cassert>
#include <random>
#include <vector>

int main() {
    std::mt19937 gen(123456u);  // Ustalony seed.
    std::uniform_int_distribution<int> count(0, 300);
    std::uniform_int_distribution<int> value(-1000, 1000);

    for (int trial = 0; trial < 1000; ++trial) {
        std::vector<int> x(count(gen));
        for (int& v : x) v = value(gen);
        auto original = x;
        std::sort(x.begin(), x.end());
        assert(std::is_sorted(x.begin(), x.end()));
        std::sort(original.begin(), original.end());
        assert(x == original);
    }
}
\end{lstlisting}

W~testach losowych przechowuj \emph{seed} i~minimalny przykład,
który ujawnił błąd. Warto również stosować \emph{shrinking}:
automatyczne zmniejszanie wejścia do możliwie prostego kontrprzykładu.

\emph{Fuzzing} automatycznie generuje wejścia, kierując się m.in.
pokryciem kodu. W~Clangu narzędzie \emph{libFuzzer} może wywoływać
poniższą funkcję z~wieloma binarnymi danymi wejściowymi.

\begin{lstlisting}[caption={Minimalny punkt wejsciowy dla libFuzzer (Clang).}]
#include <cstddef>
#include <cstdint>
#include <vector>

extern "C" int LLVMFuzzerTestOneInput(
    const std::uint8_t* data, std::size_t size) {
    // Podmien na testowany parser lub dekoder.
    std::vector<std::uint8_t> bytes(data, data + size);
    if (bytes.size() > 2 && bytes[0] == 0x42) {
        // Przypadek zalezy od danych; nie wolno zakladac ich poprawnosci.
        (void)bytes[1];
    }
    return 0;
}
\end{lstlisting}

\begin{lstlisting}[language=bash,caption={Budowanie harnessu fuzzingu w Clangu.}]
clang++ -std=c++20 -O1 -g \
  -fsanitize=fuzzer,address,undefined fuzz.cpp -o fuzz
./fuzz -max_total_time=30
\end{lstlisting}

Po zgłoszeniu awarii zapisujemy minimalny przypadek jako test
regresyjny. Fuzzer nie dowodzi braku błędów, a~jego skuteczność zależy
od jakości punktów wejściowych i~stanu pokrycia.

\section{Debugowanie: od komunikatu do przyczyny}
\label{sec:cpp18-debug}

\subsection{Informacje debugowe i rodzaje kompilacji}

Użyteczne opcje GCC/Clang to \code{-g} (symbole debugowe),
\code{-O0} (łatwe śledzenie krokowe) oraz \code{-Og} (optymalizacje
starające się zachować dobrą obserwowalność w~debuggerze). Program
kompilowany z~\code{-O2} może nie mieć osobnej zmiennej ani instrukcji
odpowiadającej każdej linii kodu źródłowego.

\begin{lstlisting}[language=bash,caption={Budowanie programu do debugowania.}]
g++ -std=c++20 -g -O0 -Wall -Wextra -Wpedantic \
    app.cpp -o app_debug

gdb ./app_debug
\end{lstlisting}

Nie wszystkie błędy są odtwarzalne przy \code{-O0}; niektóre pojawiają
się dopiero przy \code{-O2} lub \code{-O3}. Do ich analizy warto
pozostawić \code{-g} także w~wersji zoptymalizowanej i~zachować
niestripowane symbole debugowe.

\subsection{GDB: breakpoint, stos i obserwacja zmiennych}

\begin{lstlisting}[caption={Program do cwiczenia debugowania.}]
#include <stdexcept>
#include <vector>

int element_at(const std::vector<int>& v, std::size_t index) {
    return v.at(index);  // Moze zglosic std::out_of_range.
}

int main() {
    const std::vector<int> v{3, 5, 8};
    return element_at(v, 5);
}
\end{lstlisting}

\begin{lstlisting}[language=bash,caption={Typowa sesja w GDB.}]
(gdb) break element_at        # Punkt zatrzymania.
(gdb) run                     # Uruchom program.
(gdb) info args               # Argumenty funkcji.
(gdb) print index             # Wartosc zmiennej.
(gdb) backtrace               # Stos wywolan.
(gdb) next                    # Nastepna linia bez wchodzenia.
(gdb) step                    # Wejdz do wywolywanej funkcji.
(gdb) catch throw             # Zatrzymaj przy zglaszaniu wyjatku.
(gdb) continue                # Kontynuuj.
\end{lstlisting}

\code{break}, \code{watch}, \code{next}, \code{step},
\code{finish}, \code{frame}, \code{print} i~\code{bt} to podstawowe
narzędzia. \emph{Watchpoint} śledzi zmianę pamięci, ale możliwość
sprzętowego śledzenia zależy od architektury i~debuggera.

\subsection{Wątki i deadlock w debuggerze}

W~aplikacji wielowątkowej sama analiza wątku \code{main} nie
wystarczy. W~GDB:

\begin{lstlisting}[language=bash,caption={Diagnostyka wielowatkowosci w GDB.}]
(gdb) info threads
(gdb) thread 3
(gdb) bt
(gdb) thread apply all bt
(gdb) set print pretty on
\end{lstlisting}

Jeśli wątki zatrzymały się w~\code{mutex::lock()} lub
\code{condition_variable::wait()}, trzeba odtworzyć graf
zależności: który wątek posiada dany mutex i~na co czeka.
Program zakleszczony nie musi zawierać wyścigu danych, więc
ThreadSanitizer nie zastąpi analizy zakleszczeń.

\subsection{LLDB: odpowiednik na macOS i innych platformach}

LLDB oferuje analogiczne operacje:

\begin{lstlisting}[language=bash,caption={Podstawowe polecenia LLDB.}]
(lldb) breakpoint set --name element_at
(lldb) run
(lldb) frame variable
(lldb) thread backtrace
(lldb) thread list
(lldb) thread backtrace all
(lldb) next
(lldb) step
(lldb) continue
\end{lstlisting}

W~VS Code można podłączyć GDB/LLDB poprzez odpowiednie rozszerzenie,
a~następnie korzystać z~breakpointów, podglądu zmiennych i~stosów.
Dokładny adapter oraz konfiguracja \code{launch.json} zależą od
systemu, rozszerzenia i~narzędzi zainstalowanych lokalnie.

\section{Sanitizatory i narzędzia wykrywania błędów}
\label{sec:cpp18-sanitizers}

Sanitizatory dodają do programu kod kontrolujący określone klasy
błędów w~czasie wykonywania. W~odróżnieniu od debuggera mogą
wykryć nieprawidłową operację \emph{w~chwili jej wykonania}, zanim
zauważymy błędny rezultat.

\begin{center}
\small
\begin{tabularx}{\textwidth}{@{}p{0.28\textwidth}X@{}}
\toprule
Narzędzie & Główne zastosowanie\\
\midrule
AddressSanitizer (ASan) & Wyjście poza granice, use-after-free, część przypadków wycieków.\\
UndefinedBehaviorSanitizer (UBSan) & Np. przepełnienie \code{signed int}, błędne przesunięcia, null dereference, alignment.\\
ThreadSanitizer (TSan) & Wiele wyścigów danych i~część problemów synchronizacji.\\
MemorySanitizer (MSan) & Odczyty niezainicjalizowanej pamięci; wymaga specjalnych warunków instrumentacji.\\
Valgrind Memcheck & Wykrywanie wybranych błędów pamięci i~wycieków (na wspieranych systemach).\\
\bottomrule
\end{tabularx}
\end{center}

\subsection{AddressSanitizer i UndefinedBehaviorSanitizer}

Poniższy program jest \emph{celowo niepoprawny}; dostęp do elementu
po zwolnieniu pamięci jest niezdefiniowanym zachowaniem. Nie należy
interpretować jego wyniku bez instrumentacji.

\begin{lstlisting}[caption={Celowo bledny program: use-after-free.}]
int main() {
    int* p = new int(42);
    delete p;
    return *p;  // UB: odczyt z nieistniejacego obiektu.
}
\end{lstlisting}

\begin{lstlisting}[language=bash,caption={Instrumentacja ASan i UBSan; GCC/Clang.}]
g++ -std=c++20 -O1 -g -fno-omit-frame-pointer \
    -fsanitize=address,undefined -fno-sanitize-recover=all \
    bad.cpp -o bad_san
./bad_san
\end{lstlisting}

ASan i~UBSan można często połączyć. ASan/TSan/MSan nie są natomiast
zwykle łączone ze sobą w~jednym procesie: używamy osobnych
konfiguracji i~uruchomień. Raporty sanitizatorów zależą od wersji
kompilatora, wsparcia runtime i~platformy.

\subsection{ThreadSanitizer i wyścig danych}

Przykładowy wyścig występuje przy równoczesnym modyfikowaniu
zwykłego licznika. To również \emph{celowo błędny} kod:

\begin{lstlisting}[caption={Celowo bledne wspoldzielenie licznika.}]
#include <thread>

int main() {
    int x = 0;
    std::thread t1([&] { ++x; });
    std::thread t2([&] { ++x; });
    t1.join();
    t2.join();
}
\end{lstlisting}

\begin{lstlisting}[language=bash,caption={Oddzielna konfiguracja z ThreadSanitizer.}]
g++ -std=c++20 -O1 -g -pthread \
    -fsanitize=thread race.cpp -o race_tsan
./race_tsan
\end{lstlisting}

\nota{Brak zgłoszenia w~danym uruchomieniu \emph{nie} dowodzi braku
wyścigów danych. TSan ma ograniczenia platformowe, może zgłaszać
wyniki fałszywie dodatnie dla niestandardowej synchronizacji i~nie
gwarantuje wykrywania deadlocków.}

\subsection{Valgrind i testy wycieków}

Na wspieranych systemach Linux przydatna jest kontrola programów
bez kompilacyjnej instrumentacji:

\begin{lstlisting}[language=bash,caption={Przykladowe badanie pamieci przez Valgrind Memcheck.}]
valgrind --tool=memcheck --leak-check=full \
         --show-leak-kinds=all ./app_debug
\end{lstlisting}

Różne narzędzia mają inne modele wykrywania i~inne koszty
wykonania. Wybór zależy od konkretnej klasy błędów; wyniki dobrze
weryfikować na małym, odtwarzalnym przypadku testowym.

\section{Analiza statyczna, ostrzeżenia i pokrycie kodu}
\label{sec:cpp18-static}

Kompilator potrafi zgłosić wiele usterek jeszcze przed wykonaniem
programu. Warto utrzymywać ostrzeżenia w~CI i~unikać wyciszania ich
bez zrozumienia przyczyny.

\begin{lstlisting}[language=bash,caption={Przykladowe ostrzezenia GCC/Clang.}]
g++ -std=c++20 -Wall -Wextra -Wpedantic -Wconversion \
    -Wshadow -Werror -c core.cpp
\end{lstlisting}

Opcje zależą od kompilatora; \code{-Werror} bywa pożyteczne przy
stabilnej bazie kodu, lecz aktualizacja kompilatora może wprowadzić
nowe diagnostyki. \emph{clang-tidy} wykonuje zestaw analiz
statycznych, np. dotyczących nowoczesnych idiomów C++, błędów
potencjalnych i~stylu.

\begin{lstlisting}[language=bash,caption={Analiza projektu ze zbiorem komend kompilacji.}]
cmake -S . -B build -DCMAKE_EXPORT_COMPILE_COMMANDS=ON
clang-tidy -p build src/core.cpp
\end{lstlisting}

\subsection{Pokrycie testami: co liczby naprawdę oznaczają?}

\emph{Line coverage} mierzy wykonane linie, \emph{branch coverage}
wykonane gałęzie, a~\emph{function coverage} wywołane funkcje.
Wysoki procent nie oznacza sprawdzenia wszystkich kombinacji wejść
ani dowodu poprawności. Na przykład test może wykonać gałąź, nie
sprawdzając rezultatu.

\begin{lstlisting}[language=bash,caption={Pokrycie testow z GCC i gcov (przykladowy projekt).}]
g++ -std=c++20 --coverage -O0 -g core.cpp tests.cpp -o tests
./tests
gcov core.cpp
\end{lstlisting}

Dla Clanga dostępny jest osobny zestaw instrumentacji LLVM:
\begin{lstlisting}[language=bash,caption={Clang: instrumentacja i raport pokrycia.}]
clang++ -std=c++20 -O0 -g -fprofile-instr-generate \
    -fcoverage-mapping core.cpp tests.cpp -o tests
LLVM_PROFILE_FILE=tests.profraw ./tests
llvm-profdata merge -sparse tests.profraw -o tests.profdata
llvm-cov report ./tests -instr-profile=tests.profdata
\end{lstlisting}

W~szczególności należy testować \emph{warunki brzegowe}, a~nie
maksymalizować jedynie procent pokrycia.

\section{Benchmarking: jak mierzyć szybkość bez samooszukiwania}
\label{sec:cpp18-benchmark}

\subsection{Pułapki mikrobenchmarków}

Nowoczesny kompilator usuwa obliczenia o~nieobserwowalnych efektach,
wykonuje \emph{constant folding} i~wektoryzuje pętle. Dlatego mierzenie
pustej pętli lub stałego obliczenia może nie mieć związku z~realną
wydajnością. Wynik zależy ponadto od rozgrzania cache, alokacji,
obciążenia systemu, zegarów CPU, temperatury i~danych testowych.

\begin{lstlisting}[caption={Pomiar algorytmu z rzeczywistym wynikiem (bez biblioteki benchmark).}]
#include <chrono>
#include <iostream>
#include <numeric>
#include <vector>

int main() {
    std::vector<int> data(1'000'000, 1);
    using clock = std::chrono::steady_clock;
    long long checksum = 0;

    for (int repetition = 0; repetition < 10; ++repetition) {
        auto start = clock::now();
        const auto result = std::accumulate(
            data.begin(), data.end(), 0LL);
        auto finish = clock::now();
        checksum += result;  // Wynik jest wykorzystany po pomiarze.
        auto us = std::chrono::duration_cast<
            std::chrono::microseconds>(finish - start).count();
        std::cout << "Trial " << repetition << ": " << us << " us\n";
    }
    std::cout << "checksum=" << checksum << '\n';
}
\end{lstlisting}

To ilustracja, nie wzorcowy mikrobenchmark: optymalizator może
ułatwić sobie obliczenia na regularnych danych; dodatkowo
\code{std::cout} poza odcinkiem pomiarowym zmienia warunki
między próbami. Użyteczne benchmarki wymagają wielu powtórzeń,
przygotowanych danych i~analizy rozkładu czasów.

\subsection{Google Benchmark}

\emph{Google Benchmark} automatyzuje wielokrotne uruchamianie i~raportuje
miary czasowe. W~krytycznym fragmencie wykorzystujemy
\code{benchmark::DoNotOptimize}, aby wynik obliczenia nie został
trywialnie usunięty; nie wyłącza to jednak wszystkich optymalizacji
na badanym wyrażeniu.

\begin{lstlisting}[caption={Samodzielny benchmark sumowania wielu elementow.}]
// Plik: benchmark_sum.cpp
#include <benchmark/benchmark.h>
#include <numeric>
#include <vector>

static void BM_Accumulate(benchmark::State& state) {
    const std::vector<int> data(state.range(0), 1);
    for (auto _ : state) {
        const auto sum = std::accumulate(
            data.begin(), data.end(), 0LL);
        benchmark::DoNotOptimize(sum);
    }
    state.SetItemsProcessed(state.iterations() * state.range(0));
}

BENCHMARK(BM_Accumulate)->Arg(1024)->Arg(1 << 20);
BENCHMARK_MAIN();
\end{lstlisting}

\begin{lstlisting}[language={},caption={CMakeLists.txt dla przykładu Google Benchmark.}]
cmake_minimum_required(VERSION 3.20)
project(BenchmarkSum LANGUAGES CXX)
find_package(benchmark CONFIG REQUIRED)
add_executable(bench benchmark_sum.cpp)
target_compile_features(bench PRIVATE cxx_std_20)
target_link_libraries(bench PRIVATE benchmark::benchmark)
\end{lstlisting}

\begin{lstlisting}[language=bash,caption={Budowanie i powtarzanie pomiaru.}]
cmake -S . -B build -DCMAKE_BUILD_TYPE=Release
cmake --build build
./build/bench --benchmark_repetitions=10
./build/bench --benchmark_out=results.json \
              --benchmark_out_format=json
\end{lstlisting}

W~porównaniach równoległych trzeba ponadto ustalić liczbę wątków,
rozmiary bloków, warm-up, powtarzalność danych oraz ewentualne
\emph{false sharing}. Sam najkrótszy pojedynczy czas nie opisuje
stabilnej wydajności.

\section{Profilowanie CPU i pamięci}
\label{sec:cpp18-profiling}

\subsection{Profilowanie próbkowe: Linux \texttt{perf}}

Zamiast mierzyć wyłącznie całkowity czas, profiler pomaga wykrywać
\emph{hot spots}: funkcje zajmujące znaczącą część wykonania.
Przykładowe polecenia systemu Linux:

\begin{lstlisting}[language=bash,caption={Profilowanie CPU przez perf (Linux).}]
g++ -std=c++20 -O2 -g -fno-omit-frame-pointer \
    app.cpp -o app
perf stat ./app
perf record -g ./app
perf report
\end{lstlisting}

\code{perf stat} podaje liczniki zdarzeń udostępniane przez jądro
i~sprzęt, np. cykle i~instrukcje. \code{perf record} zbiera próbki
umożliwiające oglądanie gorących funkcji i~stosu. Dostępne zdarzenia
i~uprawnienia zależą od komputera i~ustawień systemu.

\subsection{Flame graph}

\emph{Flame graph} to agregat stosów wywołań: szerokość prostokąta
odpowiada udziałowi próbek, a~wysokość zagnieżdżeniu wywołań.
Nie jest to dokładna oś czasu. Przykładowy przepływ z~zewnętrznymi
narzędziami FlameGraph:

\begin{lstlisting}[language=bash,caption={Tworzenie flame graph z probek perf (Linux).}]
perf record -F 99 -g -- ./app
perf script | stackcollapse-perf.pl | flamegraph.pl > flame.svg
\end{lstlisting}

Ostatnie dwa skrypty pochodzą z~osobnego projektu FlameGraph i~muszą
być zainstalowane. W~projektach o~dużej liczbie wątków pomagają też
profile opóźnień, I/O i~blokad, a~nie tylko wykorzystania CPU.

\subsection{Alokacje i zużycie pamięci}

Wydajność pamięci to nie tylko liczba bajtów. Ważne są rozkład
alokacji, fragmentacja, lokalność cache i~czas życia obiektów.
Na wspieranych systemach można użyć m.in.:

\begin{lstlisting}[language=bash,caption={Przykladowe pomiary alokacji i pamieci.}]
valgrind --tool=massif ./app
ms_print massif.out.*
/usr/bin/time -v ./app
\end{lstlisting}

\code{massif} służy do profilowania sterty (zgodnie z~możliwościami
narzędzia); \code{/usr/bin/time -v} raportuje miary procesu, np.
\emph{maximum resident set size}, ale nie mówi, które alokacje były
odpowiedzialne. Instrumentacja pamięci również wpływa na czas pracy.

\section{Asembler i związek kodu źródłowego z wydajnością}
\label{sec:cpp18-asm}

Dwa podobne zapisy C++ mogą dać różny kod maszynowy, szczególnie po
włączeniu \code{-O2} lub \code{-O3}. Przy analizie trzeba jednak
uważać, aby nie wyciągać wniosków na podstawie sztucznych przykładów,
w których kompilator może wyeliminować całą pracę.

\begin{lstlisting}[caption={Niewielka funkcja do analizy wygenerowanego kodu.}]
#include <cstddef>

double dot(const double* a, const double* b, std::size_t n) {
    double sum = 0.0;
    for (std::size_t i = 0; i < n; ++i)
        sum += a[i] * b[i];
    return sum;
}
\end{lstlisting}

\begin{lstlisting}[language=bash,caption={Asembler i raport wektoryzacji w GCC (Linux).}]
g++ -std=c++20 -O3 -S -masm=intel dot.cpp -o dot.s

g++ -std=c++20 -O3 -fopt-info-vec-optimized \
    -c dot.cpp -o dot.o
objdump -drC dot.o
\end{lstlisting}

Dla wektoryzacji ważne są: zależności między iteracjami, wyrównanie,
\emph{aliasing} oraz semantyka arytmetyki zmiennoprzecinkowej.
Brak instrukcji SIMD nie musi oznaczać błędu kompilatora; dla
dokładnie zachowanego porządku sumowania kompilator może nie móc
zmienić kolejności operacji. Flaga \code{-ffast-math} pozwala na
istotnie inne przekształcenia numeryczne i~nie powinna być
włączana automatycznie w~programach wymagających ścisłej semantyki.

\subsection{Profile-guided optimization (PGO)}

Kompilator może wykorzystać profile zebrane na reprezentatywnych
danych, aby lepiej rozłożyć kod i~podjąć decyzje optymalizacyjne.
Schemat dla GCC:

\begin{lstlisting}[language=bash,caption={Dwa przebiegi budowania z PGO (GCC).}]
mkdir -p pgo-data
g++ -std=c++20 -O2 -fprofile-generate=pgo-data \
    app.cpp -o app
./app  # Wykonaj reprezentatywny workload.
g++ -std=c++20 -O2 -fprofile-use=pgo-data \
    app.cpp -o app
\end{lstlisting}

PGO nie zastępuje pomiarów: profil z~innego rozkładu danych może
nawet pogorszyć wyniki. Należy też zachować zgodność źródeł, flag
i~miejsc profilu pomiędzy dwoma etapami budowania.

\section{Testowanie kodu wielowątkowego}
\label{sec:cpp18-concurrency-tests}

W~systemach współbieżnych testy wynikowe muszą badać także
\emph{zakończenie pracy}, własność danych i~synchronizację.
Powtórzenie testu wiele razy zwiększa szansę wykrycia błędu, ale
nigdy nie zastępuje rozumowania o~relacjach \emph{happens-before}
(rozdziały~\ref{chap:concurrency} i~\ref{chap:cpp16-advanced-concurrency}).

\begin{lstlisting}[caption={Poprawny test licznika atomowego z jthread (C++20).}]
#include <atomic>
#include <cassert>
#include <thread>
#include <vector>

int main() {
    std::atomic<int> count{0};
    {
        std::vector<std::jthread> workers;
        for (int i = 0; i < 4; ++i) {
            workers.emplace_back([&] {
                for (int k = 0; k < 10000; ++k)
                    count.fetch_add(1, std::memory_order_relaxed);
            });
        }
    } // Destruktory jthread czekaja na wszystkie watki.
    assert(count.load(std::memory_order_relaxed) == 40000);
}
\end{lstlisting}

Do koordynowania momentu rozpoczęcia testu można wykorzystać
\code{std::latch} lub \code{std::barrier} (C++20). W~testach
\emph{timeout} zapobiega nieskończonemu oczekiwaniu w~infrastrukturze
testowej, ale nie może być uznawany za mechanizm anulowania
wątku. Zwracaj uwagę na wyścigi związane z~destrukcją fixture
oraz na kolejność zamykania wątków i~obiektów synchronizacyjnych.

\section{Praktyczny workflow: od błędu do poprawki}

Prześledźmy zalecaną kolejność działania:
\begin{enumerate}
\item \textbf{Odtwórz błąd}: zanotuj wejście, wersję programu, flagi
      kompilacji, środowisko i~opcjonalnie seed.
\item \textbf{Zminimalizuj przypadek}: mały, powtarzalny przykład jest
      łatwiejszy do analizy niż cały system.
\item \textbf{Włącz diagnostykę}: ostrzeżenia, ASan/UBSan, a~dla
      kodu współbieżnego oddzielne uruchomienie TSan.
\item \textbf{Postaw hipotezę}: zidentyfikuj zmienną lub niezmiennik,
      który naruszono; sprawdź stos i~stan programu w~debuggerze.
\item \textbf{Dodaj test regresyjny}: powinien zawodzić przed poprawką
      i~przechodzić po niej.
\item \textbf{Popraw przyczynę}: nie tylko usuń widoczny symptom.
\item \textbf{Zmierz koszt}: jeżeli poprawka ma znaczenie wydajnościowe,
      porównaj benchmarki i~profile przed oraz po zmianie.
\item \textbf{Uruchom pełen zestaw testów}: szczególnie z~różnymi
      konfiguracjami optymalizacji i~kompilatorami.
\end{enumerate}

\section{Automatyzacja w CI i macierz konfiguracji}

W~projektach utrzymywanych długo testy powinny wykonywać się
automatycznie po każdej istotnej zmianie. Praktyczna macierz obejmuje
zwykle konfigurację debugową z~sanitizatorami oraz konfigurację
optymalizowaną z~testami wynikowymi. TSan uruchamiamy w~osobnej
konfiguracji i~na wspieranym środowisku.

\begin{center}
\small
\begin{tabularx}{\textwidth}{@{}p{0.28\textwidth}X@{}}
\toprule
Konfiguracja & Główne zadanie\\
\midrule
Debug, \code{-O0 -g} & Debugowanie źródłowe i~lokalizacja błędów.\\
Debug/san, \code{-O1 -g}, ASan+UBSan & Dostęp do niepoprawnej pamięci i~wybrane UB.\\
Oddzielny build TSan & Wyścigi danych przy konkurencyjnym wykonaniu.\\
Release, \code{-O2} lub \code{-O3} & Testy produkcyjnej konfiguracji i~benchmarki.\\
Clang i GCC & Różne diagnostyki i~różne decyzje optymalizatora.\\
\bottomrule
\end{tabularx}
\end{center}

\nota{Benchmarki są wrażliwe na obciążenie współdzielonych runnerów
CI. W~zwykłym pipeline CI warto przede wszystkim weryfikować
poprawność i~duże regresje; wiarygodne mikrobenchmarki powinny
wykonywać się na kontrolowanym sprzęcie.}

\section{Najczęstsze błędy i wybór narzędzia}

\begin{center}
\small
\begin{tabularx}{\textwidth}{@{}p{0.40\textwidth}X@{}}
\toprule
Problem & Zalecane podejście\\
\midrule
Funkcja daje złe wyniki & Test jednostkowy, debugger, test regresyjny.\\
Awaria tylko dla niektórych danych & Fuzzing, test własnościowy, ASan/UBSan.\\
Podejrzenie wiszącego wskaźnika & ASan, kontrola czasu życia, debugger.\\
Podejrzenie data race & TSan, analiza synchronizacji i~\emph{happens-before}.\\
Program zawiesza się przy wielu wątkach & Stosy wszystkich wątków, graf blokad, timeout testu.\\
Wysokie zużycie pamięci & Massif, profiler alokacji, pomiar RSS.\\
Program jest wolny & \code{perf}/profiler, później benchmark gorącej funkcji.\\
Spadek szybkości po zmianie & Reprodukowalny benchmark, kontrola warunków, profil porównawczy.\\
Wątpliwości co do wektoryzacji & Raport optymalizacji i~asembler, weryfikacja numeryczna.\\
\bottomrule
\end{tabularx}
\end{center}

\section*{Podsumowanie rozdziału}
\addcontentsline{toc}{section}{Podsumowanie rozdziału}

Testy określają, czy otrzymujemy oczekiwane zachowanie; debugger
pozwala prześledzić przyczyny problemu; sanitizatory wykrywają
określone naruszenia reguł wykonywania; profiler i~benchmarki
pomagają podejmować decyzje dotyczące wydajności. Nie istnieje
jedno narzędzie, które zastępuje wszystkie pozostałe.

Praktyczna sekwencja brzmi: \emph{najpierw poprawność, potem
bezpieczeństwo i~reprodukcyjność, następnie pomiary, wreszcie
optymalizacja}. Każda wykryta usterka powinna prowadzić do testu
regresyjnego, a~każda deklarowana poprawa wydajności --- do
porównywalnego pomiaru. Dla systemów wielowątkowych konieczne jest
ponadto świadome rozumienie synchronizacji, czasu życia obiektów
oraz gwarancji modelu pamięci C++.

\section*{Literatura i dokumentacja}
\addcontentsline{toc}{section}{Literatura i dokumentacja}
\begingroup
\sloppy
\begin{itemize}
\item \url{https://google.github.io/googletest/primer.html} ---
      GoogleTest: asercje, fixture i organizacja testów.
\item \url{https://github.com/google/googletest/blob/main/docs/quickstart-cmake.md}
      --- integracja GoogleTest i CMake.
\item \url{https://cmake.org/cmake/help/latest/module/GoogleTest.html}
      --- \code{gtest_discover_tests} i CTest.
\item \url{https://github.com/google/benchmark/blob/main/docs/user_guide.md}
      --- Google Benchmark i stabilne pomiary.
\item \url{https://clang.llvm.org/docs/AddressSanitizer.html} --- ASan.
\item \url{https://clang.llvm.org/docs/UndefinedBehaviorSanitizer.html}
      --- UBSan.
\item \url{https://clang.llvm.org/docs/ThreadSanitizer.html} --- TSan.
\item \url{https://clang.llvm.org/docs/LibFuzzer.html} --- fuzzing.
\item \url{https://sourceware.org/gdb/current/onlinedocs/gdb.html} --- GDB.
\item \url{https://lldb.llvm.org/use/tutorial.html} --- LLDB.
\item \url{https://valgrind.org/docs/manual/manual.html} --- Memcheck i Massif.
\item \url{https://perf.wiki.kernel.org/} --- profiler \code{perf}.
\item Brendan Gregg, \emph{Systems Performance}, wyd. 2,
      Addison-Wesley, 2020.
\end{itemize}
\endgroup


\chapter{Testowanie, debugowanie i profilowanie C++}
\label{chap:cpp18-test-debug-profile}

\section{Trzy różne pytania o program}

Po napisaniu i skompilowaniu programu nie wiemy jeszcze, czy jest on
poprawny, bezpieczny i wydajny. Każdą z tych własności sprawdza się
innymi metodami:
\begin{description}[style=nextline]
\item[Testowanie (\emph{testing})]
Czy program daje oczekiwane wyniki na wybranych danych oraz czy
spełnia wymagane kontrakty? Testy wykrywają błędy, lecz ich pozytywny
wynik nie jest ogólnym dowodem poprawności.
\item[Debugowanie (\emph{debugging})]
Gdzie i dlaczego program odbiega od oczekiwanego zachowania?
Debugger pozwala obserwować stos wywołań, zmienne, wątki i kolejne
instrukcje.
\item[Profilowanie (\emph{profiling})]
Gdzie program zużywa czas procesora, pamięć i inne zasoby?
Profilowanie wskazuje rzeczywiste wąskie gardła; samo zgadywanie
na podstawie długości kodu często prowadzi do błędnych optymalizacji.
\end{description}

Istotne jest także rozróżnienie \emph{testu poprawności} od
\emph{benchmarku}: pierwszy sprawdza warunek logiczny, drugi mierzy
wielkość fizyczną (np. czas, przepustowość, alokacje). Szybki, lecz
błędny algorytm nie jest dobrym rozwiązaniem.

\nota{W~tym rozdziale podstawowym standardem jest C++20. GoogleTest,
Google Benchmark, GDB, LLDB, Valgrind i~narzędzia \code{perf} nie są
częścią standardowej biblioteki C++; wymagają odpowiednich instalacji.
Polecenia zależą od systemu operacyjnego i~kompilatora.}

\section{Strategia testowania i testowalny projekt}
\label{sec:cpp18-strategy}

\subsection{Poziomy testów}

\begin{center}
\small
\begin{tabularx}{\textwidth}{@{}p{0.29\textwidth}X@{}}
\toprule
Rodzaj testu & Co sprawdza i kiedy jest przydatny?\\
\midrule
Jednostkowy (\emph{unit}) & Mała, izolowana funkcja lub klasa; szybka informacja zwrotna.\\
Integracyjny & Współdziałanie komponentów, np. parsera i~bazy danych.\\
Systemowy / end-to-end & Cały scenariusz użytkownika, często z~rzeczywistym wejściem/wyjściem.\\
Regresyjny & Przypadek, który wcześniej ujawniał błąd; nie powinien powrócić.\\
Własnościowy / losowy & Niezmienniki dla wielu zestawów danych, nie tylko pojedyncze wyniki.\\
Wydajnościowy & Czas, przepustowość, opóźnienie, pamięć; wymaga kontrolowanego środowiska.\\
\bottomrule
\end{tabularx}
\end{center}

Dobry test powinien być \emph{powtarzalny}, \emph{niezależny} od innych
testów i~możliwie mały. W~praktyce wydzielamy obliczenia od efektów
ubocznych (sieci, zegara, plików, losowości), dzięki czemu zależności
możemy zastąpić kontrolowanymi implementacjami.

\subsection{Kontrakty, przypadki brzegowe i błędy numeryczne}

Rozważmy funkcję obliczającą średnią. Jej kontrakt powinien ustalać
zachowanie dla pustego zakresu; poniżej zgłaszamy wyjątek.

\begin{lstlisting}[caption={Prosta funkcja z jawnym kontraktem dla pustego zakresu (C++20).}]
#include <numeric>
#include <span>
#include <stdexcept>

double average(std::span<const double> xs) {
    if (xs.empty())
        throw std::invalid_argument("empty range");
    return std::accumulate(xs.begin(), xs.end(), 0.0)
         / static_cast<double>(xs.size());
}
\end{lstlisting}

Testujemy nie tylko wartości typowe, lecz także rozmiary $0$ i~$1$,
wartości ujemne, powtórzenia, skrajne zakresy oraz dane naruszające
warunki wejściowe. Dla liczb zmiennoprzecinkowych zwykła równość
\code{a == b} nie zawsze jest właściwa: wynik może zależeć od porządku
dodawania i~optymalizacji. Stosujemy tolerancję dostosowaną do skali
i~wymagań zadania, a~nie arbitralne \code{1e-6} dla każdego algorytmu.

\section{Testy bez zewnętrznego frameworka}

\subsection{\texttt{assert} i \texttt{static\_assert}}

\code{assert(expr)} z~nagłówka \code{<cassert>} sprawdza warunek
w~czasie działania, lecz może zostać usunięte przy zdefiniowanym
\code{NDEBUG}. \code{static_assert} sprawdza warunek w~czasie
kompilacji. W~szczególności zwykłe \code{assert} nie nadaje się do
walidacji niepoprawnych danych wejściowych w~kodzie produkcyjnym.

\begin{lstlisting}[caption={Asercja kompilacyjna i kontrola zachowania funkcji.}]
#include <cassert>
#include <concepts>
#include <type_traits>

constexpr int square(int x) { return x * x; }
static_assert(square(6) == 36);
static_assert(std::integral<int>);
static_assert(!std::is_copy_constructible_v<void>);

int main() {
    assert(square(7) == 49);
}
\end{lstlisting}

\nota{Przy \code{-DNDEBUG} argument \code{assert} może w~ogóle nie być
wykonywany. Nie wolno umieszczać w~nim niezbędnych efektów ubocznych,
np.~\code{assert(++counter == 1)}. \code{static_assert} również nie
zastępuje testów dla danych znanych dopiero podczas wykonania.}

\subsection{Niezależne testy danych wejściowych}

Poniższą funkcję \code{main} umieść pod definicją \code{average}
z~poprzedniego listingu. Test nie wymaga zewnętrznej biblioteki.
Sprawdza wynik, wyjątek i~przypadek jednoelementowy. Kompiluj go
bez \code{-DNDEBUG}, aby wyrażenia \code{assert} były wykonywane.

\begin{lstlisting}[caption={Test funkcji average bez frameworka; dołącz pod jej definicją (C++20).}]
#include <cassert>
#include <cmath>
#include <stdexcept>
#include <vector>

// Zakladamy definicje average z poprzedniego przykladu.
int main() {
    const std::vector<double> xs{1.0, 2.0, 3.0, 4.0};
    assert(std::abs(average(xs) - 2.5) < 1e-12);
    const std::vector<double> one{7.0};
    assert(average(one) == 7.0);

    bool caught = false;
    const std::vector<double> empty;
    try {
        average(empty);
    } catch (const std::invalid_argument&) {
        caught = true;
    }
    assert(caught);
}
\end{lstlisting}

W~realnym projekcie lepiej użyć frameworka, który rejestruje testy,
raportuje niepowodzenia i~pozwala uruchamiać wybrane przypadki.

\section{GoogleTest: testy jednostkowe w praktyce}
\label{sec:cpp18-gtest}

Biblioteka \emph{GoogleTest} (często \code{gtest}) udostępnia makra
\code{TEST}, \code{TEST_F}, \code{EXPECT_*} i~\code{ASSERT_*}.
Assercje \code{EXPECT_*} zapisują niepowodzenie i~pozwalają wykonać
resztę funkcji testowej, podczas gdy \code{ASSERT_*} przerywają
bieżącą funkcję testową. To nie są mechanizmy obsługi wyjątków C++.

\begin{lstlisting}[caption={GoogleTest: wyniki, wyjatki i przyblizona rownosc.}]
// Plik: average_test.cpp
#include <gtest/gtest.h>
#include <stdexcept>
#include <vector>

#include "average.hpp"
// Funkcja average jest zdefiniowana w pliku average.cpp.
TEST(AverageTest, ComputesMean) {
    const std::vector<double> xs{2.0, 4.0, 6.0};
    EXPECT_NEAR(average(xs), 4.0, 1e-12);
}

TEST(AverageTest, RejectsEmptyInput) {
    const std::vector<double> empty;
    EXPECT_THROW(average(empty), std::invalid_argument);
}

TEST(AverageTest, HandlesSingleton) {
    const std::vector<double> xs{7.0};
    ASSERT_EQ(xs.size(), 1u);
    EXPECT_DOUBLE_EQ(average(xs), 7.0);
}
\end{lstlisting}

\code{EXPECT_NEAR(a,b,tol)} sprawdza bezwzględną różnicę, dlatego
niekiedy lepiej użyć własnego testu tolerancji względnej. Makro
\code{EXPECT_DOUBLE_EQ} porównuje liczby z~tolerancją opartą na ULP;
nie należy utożsamiać go z~matematyczną równością liczb rzeczywistych.

\subsection{Fixture i testy parametryzowane}

\emph{Fixture} przygotowuje stan wspólny dla testów. Każdy test
\code{TEST_F} otrzymuje osobny egzemplarz fixture, dzięki czemu zmiany
jednego testu nie powinny wpływać na inny.

\begin{lstlisting}[caption={Fixture z przygotowanym wektorem.}]
#include <gtest/gtest.h>
#include <algorithm>
#include <vector>

class VectorTest : public ::testing::Test {
protected:
    std::vector<int> values{3, 1, 4, 1, 5};
};

TEST_F(VectorTest, SortsAscending) {
    std::sort(values.begin(), values.end());
    EXPECT_TRUE(std::is_sorted(values.begin(), values.end()));
}

TEST_F(VectorTest, StartsFromOriginalState) {
    EXPECT_EQ(values.front(), 3);
    EXPECT_EQ(values.size(), 5u);
}
\end{lstlisting}

Dla wielu par wejście--wynik przydatne są testy parametryzowane:

\begin{lstlisting}[caption={GoogleTest: test parametryzowany z zestawem wartosci.}]
#include <gtest/gtest.h>
#include <tuple>

constexpr int square(int x) { return x * x; }

class SquareTest
    : public ::testing::TestWithParam<std::tuple<int, int>> {};

TEST_P(SquareTest, MatchesExpected) {
    auto [x, expected] = GetParam();
    EXPECT_EQ(square(x), expected);
}

INSTANTIATE_TEST_SUITE_P(
    Values, SquareTest,
    ::testing::Values(
        std::tuple{0, 0},
        std::tuple{2, 4},
        std::tuple{-3, 9}));
\end{lstlisting}

\subsection{Mocki i testowanie współpracy komponentów}

\emph{GoogleMock} pozwala tworzyć kontrolowane implementacje
interfejsów i~sprawdzać wywołania między obiektami. Używaj go tam,
gdzie zachowanie naprawdę zależy od interakcji; nadmierne mockowanie
utrwala przypadkowe szczegóły implementacji.

\begin{lstlisting}[caption={GoogleMock: sprawdzenie pojedynczego wywolania.}]
#include <gmock/gmock.h>
#include <string_view>

struct Sender {
    virtual ~Sender() = default;
    virtual bool send(std::string_view) = 0;
};

struct MockSender : Sender {
    MOCK_METHOD(bool, send, (std::string_view), (override));
};

void notify(Sender& s) { (void)s.send("ready"); }

TEST(NotificationTest, SendsOneMessage) {
    MockSender sender;
    EXPECT_CALL(sender, send("ready"))
        .Times(1).WillOnce(::testing::Return(true));
    notify(sender);
}
\end{lstlisting}

Własne wyjątki, klasy niewspierające kopiowania i~asynchroniczne
wywołania wymagają odpowiedniego projektu fixture/mocków.
Dla operacji wykonywanych w~innym wątku szczególnie istotne jest
kontrolowane oczekiwanie na koniec pracy przed zniszczeniem mocka.

\subsection{CMake i CTest: mały projekt wieloplikowy}
\label{sec:cpp18-cmake-test}

Poniższa organizacja zakłada, że \code{average} ma deklarację w~pliku
\code{average.hpp}, definicję w~\code{average.cpp}, a~testy w~pliku
\code{average_test.cpp}. W~przykładzie \code{average_test.cpp}
należy dodać \verb|#include "average.hpp"| zamiast deklarować
funkcję ręcznie. Plik nagłówkowy może wyglądać następująco:

\begin{lstlisting}[caption={Plik average.hpp: deklaracja API testowanej funkcji.}]
#pragma once
#include <span>
double average(std::span<const double> xs);
\end{lstlisting}

Plik \code{average.cpp} zawiera kod pierwszego przykładu
z~tej części rozdziału (z~dołączonym \code{average.hpp});
w~pliku testów umieszczamy testy z~sekcji~\ref{sec:cpp18-gtest}.

\begin{lstlisting}[language=bash,caption={Struktura projektu z GoogleTest.}]
project/
  CMakeLists.txt
  average.hpp
  average.cpp
  average_test.cpp
\end{lstlisting}

\begin{lstlisting}[language=bash,caption={CMakeLists.txt: biblioteka, testy i automatyczna rejestracja w CTest.}]
cmake_minimum_required(VERSION 3.20)
project(cpp18_tests LANGUAGES CXX)

add_library(core average.cpp)
target_compile_features(core PUBLIC cxx_std_20)
target_include_directories(core PUBLIC ${CMAKE_CURRENT_SOURCE_DIR})

# GoogleTest musi byc zainstalowany lub dostarczony przez menedzer pakietow.
find_package(GTest CONFIG REQUIRED)

add_executable(core_tests average_test.cpp)
target_link_libraries(core_tests PRIVATE core GTest::gtest_main)

enable_testing()
include(GoogleTest)
gtest_discover_tests(core_tests)
\end{lstlisting}

\begin{lstlisting}[language=bash,caption={Budowanie i uruchamianie testow przez CMake/CTest.}]
cmake -S . -B build -DCMAKE_BUILD_TYPE=Debug
cmake --build build -j
ctest --test-dir build --output-on-failure
./build/core_tests --gtest_filter='AverageTest.*'
\end{lstlisting}

\nota{W~systemach z~generatorem wielokonfiguracyjnym (np.~Visual
Studio) wariant budowania wybiera się przez \code{--config Debug}.
Przykład wymaga zainstalowanego GoogleTest; alternatywą jest
\code{FetchContent} z~\emph{przypiętą} wersją zależności, a~nie
pobieranie nieokreślonej wersji podczas każdego budowania.}

\section{Testy własnościowe, generatywne i fuzzing}
\label{sec:cpp18-property}

Sprawdzanie setek wartości jest bardziej użyteczne, jeśli testuje
\emph{własność}, a~nie jedynie listę ręcznie wpisanych wyników.
Na przykład po posortowaniu ciągu musi on być niemalejący i~mieć
tę samą multizbiorowość wartości co przed sortowaniem.

\begin{lstlisting}[caption={Powtarzalny test wlasnosciowy przy uzyciu generatora losowego.}]
#include <algorithm>
#include <cassert>
#include <random>
#include <vector>

int main() {
    std::mt19937 gen(123456u);  // Ustalony seed.
    std::uniform_int_distribution<int> count(0, 300);
    std::uniform_int_distribution<int> value(-1000, 1000);

    for (int trial = 0; trial < 1000; ++trial) {
        std::vector<int> x(count(gen));
        for (int& v : x) v = value(gen);
        auto original = x;
        std::sort(x.begin(), x.end());
        assert(std::is_sorted(x.begin(), x.end()));
        std::sort(original.begin(), original.end());
        assert(x == original);
    }
}
\end{lstlisting}

W~testach losowych przechowuj \emph{seed} i~minimalny przykład,
który ujawnił błąd. Warto również stosować \emph{shrinking}:
automatyczne zmniejszanie wejścia do możliwie prostego kontrprzykładu.

\emph{Fuzzing} automatycznie generuje wejścia, kierując się m.in.
pokryciem kodu. W~Clangu narzędzie \emph{libFuzzer} może wywoływać
poniższą funkcję z~wieloma binarnymi danymi wejściowymi.

\begin{lstlisting}[caption={Minimalny punkt wejsciowy dla libFuzzer (Clang).}]
#include <cstddef>
#include <cstdint>
#include <vector>

extern "C" int LLVMFuzzerTestOneInput(
    const std::uint8_t* data, std::size_t size) {
    // Podmien na testowany parser lub dekoder.
    std::vector<std::uint8_t> bytes(data, data + size);
    if (bytes.size() > 2 && bytes[0] == 0x42) {
        // Przypadek zalezy od danych; nie wolno zakladac ich poprawnosci.
        (void)bytes[1];
    }
    return 0;
}
\end{lstlisting}

\begin{lstlisting}[language=bash,caption={Budowanie harnessu fuzzingu w Clangu.}]
clang++ -std=c++20 -O1 -g \
  -fsanitize=fuzzer,address,undefined fuzz.cpp -o fuzz
./fuzz -max_total_time=30
\end{lstlisting}

Po zgłoszeniu awarii zapisujemy minimalny przypadek jako test
regresyjny. Fuzzer nie dowodzi braku błędów, a~jego skuteczność zależy
od jakości punktów wejściowych i~stanu pokrycia.

\section{Debugowanie: od komunikatu do przyczyny}
\label{sec:cpp18-debug}

\subsection{Informacje debugowe i rodzaje kompilacji}

Użyteczne opcje GCC/Clang to \code{-g} (symbole debugowe),
\code{-O0} (łatwe śledzenie krokowe) oraz \code{-Og} (optymalizacje
starające się zachować dobrą obserwowalność w~debuggerze). Program
kompilowany z~\code{-O2} może nie mieć osobnej zmiennej ani instrukcji
odpowiadającej każdej linii kodu źródłowego.

\begin{lstlisting}[language=bash,caption={Budowanie programu do debugowania.}]
g++ -std=c++20 -g -O0 -Wall -Wextra -Wpedantic \
    app.cpp -o app_debug

gdb ./app_debug
\end{lstlisting}

Nie wszystkie błędy są odtwarzalne przy \code{-O0}; niektóre pojawiają
się dopiero przy \code{-O2} lub \code{-O3}. Do ich analizy warto
pozostawić \code{-g} także w~wersji zoptymalizowanej i~zachować
niestripowane symbole debugowe.

\subsection{GDB: breakpoint, stos i obserwacja zmiennych}

\begin{lstlisting}[caption={Program do cwiczenia debugowania.}]
#include <stdexcept>
#include <vector>

int element_at(const std::vector<int>& v, std::size_t index) {
    return v.at(index);  // Moze zglosic std::out_of_range.
}

int main() {
    const std::vector<int> v{3, 5, 8};
    return element_at(v, 5);
}
\end{lstlisting}

\begin{lstlisting}[language=bash,caption={Typowa sesja w GDB.}]
(gdb) break element_at        # Punkt zatrzymania.
(gdb) run                     # Uruchom program.
(gdb) info args               # Argumenty funkcji.
(gdb) print index             # Wartosc zmiennej.
(gdb) backtrace               # Stos wywolan.
(gdb) next                    # Nastepna linia bez wchodzenia.
(gdb) step                    # Wejdz do wywolywanej funkcji.
(gdb) catch throw             # Zatrzymaj przy zglaszaniu wyjatku.
(gdb) continue                # Kontynuuj.
\end{lstlisting}

\code{break}, \code{watch}, \code{next}, \code{step},
\code{finish}, \code{frame}, \code{print} i~\code{bt} to podstawowe
narzędzia. \emph{Watchpoint} śledzi zmianę pamięci, ale możliwość
sprzętowego śledzenia zależy od architektury i~debuggera.

\subsection{Wątki i deadlock w debuggerze}

W~aplikacji wielowątkowej sama analiza wątku \code{main} nie
wystarczy. W~GDB:

\begin{lstlisting}[language=bash,caption={Diagnostyka wielowatkowosci w GDB.}]
(gdb) info threads
(gdb) thread 3
(gdb) bt
(gdb) thread apply all bt
(gdb) set print pretty on
\end{lstlisting}

Jeśli wątki zatrzymały się w~\code{mutex::lock()} lub
\code{condition_variable::wait()}, trzeba odtworzyć graf
zależności: który wątek posiada dany mutex i~na co czeka.
Program zakleszczony nie musi zawierać wyścigu danych, więc
ThreadSanitizer nie zastąpi analizy zakleszczeń.

\subsection{LLDB: odpowiednik na macOS i innych platformach}

LLDB oferuje analogiczne operacje:

\begin{lstlisting}[language=bash,caption={Podstawowe polecenia LLDB.}]
(lldb) breakpoint set --name element_at
(lldb) run
(lldb) frame variable
(lldb) thread backtrace
(lldb) thread list
(lldb) thread backtrace all
(lldb) next
(lldb) step
(lldb) continue
\end{lstlisting}

W~VS Code można podłączyć GDB/LLDB poprzez odpowiednie rozszerzenie,
a~następnie korzystać z~breakpointów, podglądu zmiennych i~stosów.
Dokładny adapter oraz konfiguracja \code{launch.json} zależą od
systemu, rozszerzenia i~narzędzi zainstalowanych lokalnie.

\section{Sanitizatory i narzędzia wykrywania błędów}
\label{sec:cpp18-sanitizers}

Sanitizatory dodają do programu kod kontrolujący określone klasy
błędów w~czasie wykonywania. W~odróżnieniu od debuggera mogą
wykryć nieprawidłową operację \emph{w~chwili jej wykonania}, zanim
zauważymy błędny rezultat.

\begin{center}
\small
\begin{tabularx}{\textwidth}{@{}p{0.28\textwidth}X@{}}
\toprule
Narzędzie & Główne zastosowanie\\
\midrule
AddressSanitizer (ASan) & Wyjście poza granice, use-after-free, część przypadków wycieków.\\
UndefinedBehaviorSanitizer (UBSan) & Np. przepełnienie \code{signed int}, błędne przesunięcia, null dereference, alignment.\\
ThreadSanitizer (TSan) & Wiele wyścigów danych i~część problemów synchronizacji.\\
MemorySanitizer (MSan) & Odczyty niezainicjalizowanej pamięci; wymaga specjalnych warunków instrumentacji.\\
Valgrind Memcheck & Wykrywanie wybranych błędów pamięci i~wycieków (na wspieranych systemach).\\
\bottomrule
\end{tabularx}
\end{center}

\subsection{AddressSanitizer i UndefinedBehaviorSanitizer}

Poniższy program jest \emph{celowo niepoprawny}; dostęp do elementu
po zwolnieniu pamięci jest niezdefiniowanym zachowaniem. Nie należy
interpretować jego wyniku bez instrumentacji.

\begin{lstlisting}[caption={Celowo bledny program: use-after-free.}]
int main() {
    int* p = new int(42);
    delete p;
    return *p;  // UB: odczyt z nieistniejacego obiektu.
}
\end{lstlisting}

\begin{lstlisting}[language=bash,caption={Instrumentacja ASan i UBSan; GCC/Clang.}]
g++ -std=c++20 -O1 -g -fno-omit-frame-pointer \
    -fsanitize=address,undefined -fno-sanitize-recover=all \
    bad.cpp -o bad_san
./bad_san
\end{lstlisting}

ASan i~UBSan można często połączyć. ASan/TSan/MSan nie są natomiast
zwykle łączone ze sobą w~jednym procesie: używamy osobnych
konfiguracji i~uruchomień. Raporty sanitizatorów zależą od wersji
kompilatora, wsparcia runtime i~platformy.

\subsection{ThreadSanitizer i wyścig danych}

Przykładowy wyścig występuje przy równoczesnym modyfikowaniu
zwykłego licznika. To również \emph{celowo błędny} kod:

\begin{lstlisting}[caption={Celowo bledne wspoldzielenie licznika.}]
#include <thread>

int main() {
    int x = 0;
    std::thread t1([&] { ++x; });
    std::thread t2([&] { ++x; });
    t1.join();
    t2.join();
}
\end{lstlisting}

\begin{lstlisting}[language=bash,caption={Oddzielna konfiguracja z ThreadSanitizer.}]
g++ -std=c++20 -O1 -g -pthread \
    -fsanitize=thread race.cpp -o race_tsan
./race_tsan
\end{lstlisting}

\nota{Brak zgłoszenia w~danym uruchomieniu \emph{nie} dowodzi braku
wyścigów danych. TSan ma ograniczenia platformowe, może zgłaszać
wyniki fałszywie dodatnie dla niestandardowej synchronizacji i~nie
gwarantuje wykrywania deadlocków.}

\subsection{Valgrind i testy wycieków}

Na wspieranych systemach Linux przydatna jest kontrola programów
bez kompilacyjnej instrumentacji:

\begin{lstlisting}[language=bash,caption={Przykladowe badanie pamieci przez Valgrind Memcheck.}]
valgrind --tool=memcheck --leak-check=full \
         --show-leak-kinds=all ./app_debug
\end{lstlisting}

Różne narzędzia mają inne modele wykrywania i~inne koszty
wykonania. Wybór zależy od konkretnej klasy błędów; wyniki dobrze
weryfikować na małym, odtwarzalnym przypadku testowym.

\section{Analiza statyczna, ostrzeżenia i pokrycie kodu}
\label{sec:cpp18-static}

Kompilator potrafi zgłosić wiele usterek jeszcze przed wykonaniem
programu. Warto utrzymywać ostrzeżenia w~CI i~unikać wyciszania ich
bez zrozumienia przyczyny.

\begin{lstlisting}[language=bash,caption={Przykladowe ostrzezenia GCC/Clang.}]
g++ -std=c++20 -Wall -Wextra -Wpedantic -Wconversion \
    -Wshadow -Werror -c core.cpp
\end{lstlisting}

Opcje zależą od kompilatora; \code{-Werror} bywa pożyteczne przy
stabilnej bazie kodu, lecz aktualizacja kompilatora może wprowadzić
nowe diagnostyki. \emph{clang-tidy} wykonuje zestaw analiz
statycznych, np. dotyczących nowoczesnych idiomów C++, błędów
potencjalnych i~stylu.

\begin{lstlisting}[language=bash,caption={Analiza projektu ze zbiorem komend kompilacji.}]
cmake -S . -B build -DCMAKE_EXPORT_COMPILE_COMMANDS=ON
clang-tidy -p build src/core.cpp
\end{lstlisting}

\subsection{Pokrycie testami: co liczby naprawdę oznaczają?}

\emph{Line coverage} mierzy wykonane linie, \emph{branch coverage}
wykonane gałęzie, a~\emph{function coverage} wywołane funkcje.
Wysoki procent nie oznacza sprawdzenia wszystkich kombinacji wejść
ani dowodu poprawności. Na przykład test może wykonać gałąź, nie
sprawdzając rezultatu.

\begin{lstlisting}[language=bash,caption={Pokrycie testow z GCC i gcov (przykladowy projekt).}]
g++ -std=c++20 --coverage -O0 -g core.cpp tests.cpp -o tests
./tests
gcov core.cpp
\end{lstlisting}

Dla Clanga dostępny jest osobny zestaw instrumentacji LLVM:
\begin{lstlisting}[language=bash,caption={Clang: instrumentacja i raport pokrycia.}]
clang++ -std=c++20 -O0 -g -fprofile-instr-generate \
    -fcoverage-mapping core.cpp tests.cpp -o tests
LLVM_PROFILE_FILE=tests.profraw ./tests
llvm-profdata merge -sparse tests.profraw -o tests.profdata
llvm-cov report ./tests -instr-profile=tests.profdata
\end{lstlisting}

W~szczególności należy testować \emph{warunki brzegowe}, a~nie
maksymalizować jedynie procent pokrycia.

\section{Benchmarking: jak mierzyć szybkość bez samooszukiwania}
\label{sec:cpp18-benchmark}

\subsection{Pułapki mikrobenchmarków}

Nowoczesny kompilator usuwa obliczenia o~nieobserwowalnych efektach,
wykonuje \emph{constant folding} i~wektoryzuje pętle. Dlatego mierzenie
pustej pętli lub stałego obliczenia może nie mieć związku z~realną
wydajnością. Wynik zależy ponadto od rozgrzania cache, alokacji,
obciążenia systemu, zegarów CPU, temperatury i~danych testowych.

\begin{lstlisting}[caption={Pomiar algorytmu z rzeczywistym wynikiem (bez biblioteki benchmark).}]
#include <chrono>
#include <iostream>
#include <numeric>
#include <vector>

int main() {
    std::vector<int> data(1'000'000, 1);
    using clock = std::chrono::steady_clock;
    long long checksum = 0;

    for (int repetition = 0; repetition < 10; ++repetition) {
        auto start = clock::now();
        const auto result = std::accumulate(
            data.begin(), data.end(), 0LL);
        auto finish = clock::now();
        checksum += result;  // Wynik jest wykorzystany po pomiarze.
        auto us = std::chrono::duration_cast<
            std::chrono::microseconds>(finish - start).count();
        std::cout << "Trial " << repetition << ": " << us << " us\n";
    }
    std::cout << "checksum=" << checksum << '\n';
}
\end{lstlisting}

To ilustracja, nie wzorcowy mikrobenchmark: optymalizator może
ułatwić sobie obliczenia na regularnych danych; dodatkowo
\code{std::cout} poza odcinkiem pomiarowym zmienia warunki
między próbami. Użyteczne benchmarki wymagają wielu powtórzeń,
przygotowanych danych i~analizy rozkładu czasów.

\subsection{Google Benchmark}

\emph{Google Benchmark} automatyzuje wielokrotne uruchamianie i~raportuje
miary czasowe. W~krytycznym fragmencie wykorzystujemy
\code{benchmark::DoNotOptimize}, aby wynik obliczenia nie został
trywialnie usunięty; nie wyłącza to jednak wszystkich optymalizacji
na badanym wyrażeniu.

\begin{lstlisting}[caption={Samodzielny benchmark sumowania wielu elementow.}]
// Plik: benchmark_sum.cpp
#include <benchmark/benchmark.h>
#include <numeric>
#include <vector>

static void BM_Accumulate(benchmark::State& state) {
    const std::vector<int> data(state.range(0), 1);
    for (auto _ : state) {
        const auto sum = std::accumulate(
            data.begin(), data.end(), 0LL);
        benchmark::DoNotOptimize(sum);
    }
    state.SetItemsProcessed(state.iterations() * state.range(0));
}

BENCHMARK(BM_Accumulate)->Arg(1024)->Arg(1 << 20);
BENCHMARK_MAIN();
\end{lstlisting}

\begin{lstlisting}[language={},caption={CMakeLists.txt dla przykładu Google Benchmark.}]
cmake_minimum_required(VERSION 3.20)
project(BenchmarkSum LANGUAGES CXX)
find_package(benchmark CONFIG REQUIRED)
add_executable(bench benchmark_sum.cpp)
target_compile_features(bench PRIVATE cxx_std_20)
target_link_libraries(bench PRIVATE benchmark::benchmark)
\end{lstlisting}

\begin{lstlisting}[language=bash,caption={Budowanie i powtarzanie pomiaru.}]
cmake -S . -B build -DCMAKE_BUILD_TYPE=Release
cmake --build build
./build/bench --benchmark_repetitions=10
./build/bench --benchmark_out=results.json \
              --benchmark_out_format=json
\end{lstlisting}

W~porównaniach równoległych trzeba ponadto ustalić liczbę wątków,
rozmiary bloków, warm-up, powtarzalność danych oraz ewentualne
\emph{false sharing}. Sam najkrótszy pojedynczy czas nie opisuje
stabilnej wydajności.

\section{Profilowanie CPU i pamięci}
\label{sec:cpp18-profiling}

\subsection{Profilowanie próbkowe: Linux \texttt{perf}}

Zamiast mierzyć wyłącznie całkowity czas, profiler pomaga wykrywać
\emph{hot spots}: funkcje zajmujące znaczącą część wykonania.
Przykładowe polecenia systemu Linux:

\begin{lstlisting}[language=bash,caption={Profilowanie CPU przez perf (Linux).}]
g++ -std=c++20 -O2 -g -fno-omit-frame-pointer \
    app.cpp -o app
perf stat ./app
perf record -g ./app
perf report
\end{lstlisting}

\code{perf stat} podaje liczniki zdarzeń udostępniane przez jądro
i~sprzęt, np. cykle i~instrukcje. \code{perf record} zbiera próbki
umożliwiające oglądanie gorących funkcji i~stosu. Dostępne zdarzenia
i~uprawnienia zależą od komputera i~ustawień systemu.

\subsection{Flame graph}

\emph{Flame graph} to agregat stosów wywołań: szerokość prostokąta
odpowiada udziałowi próbek, a~wysokość zagnieżdżeniu wywołań.
Nie jest to dokładna oś czasu. Przykładowy przepływ z~zewnętrznymi
narzędziami FlameGraph:

\begin{lstlisting}[language=bash,caption={Tworzenie flame graph z probek perf (Linux).}]
perf record -F 99 -g -- ./app
perf script | stackcollapse-perf.pl | flamegraph.pl > flame.svg
\end{lstlisting}

Ostatnie dwa skrypty pochodzą z~osobnego projektu FlameGraph i~muszą
być zainstalowane. W~projektach o~dużej liczbie wątków pomagają też
profile opóźnień, I/O i~blokad, a~nie tylko wykorzystania CPU.

\subsection{Alokacje i zużycie pamięci}

Wydajność pamięci to nie tylko liczba bajtów. Ważne są rozkład
alokacji, fragmentacja, lokalność cache i~czas życia obiektów.
Na wspieranych systemach można użyć m.in.:

\begin{lstlisting}[language=bash,caption={Przykladowe pomiary alokacji i pamieci.}]
valgrind --tool=massif ./app
ms_print massif.out.*
/usr/bin/time -v ./app
\end{lstlisting}

\code{massif} służy do profilowania sterty (zgodnie z~możliwościami
narzędzia); \code{/usr/bin/time -v} raportuje miary procesu, np.
\emph{maximum resident set size}, ale nie mówi, które alokacje były
odpowiedzialne. Instrumentacja pamięci również wpływa na czas pracy.

\section{Asembler i związek kodu źródłowego z wydajnością}
\label{sec:cpp18-asm}

Dwa podobne zapisy C++ mogą dać różny kod maszynowy, szczególnie po
włączeniu \code{-O2} lub \code{-O3}. Przy analizie trzeba jednak
uważać, aby nie wyciągać wniosków na podstawie sztucznych przykładów,
w których kompilator może wyeliminować całą pracę.

\begin{lstlisting}[caption={Niewielka funkcja do analizy wygenerowanego kodu.}]
#include <cstddef>

double dot(const double* a, const double* b, std::size_t n) {
    double sum = 0.0;
    for (std::size_t i = 0; i < n; ++i)
        sum += a[i] * b[i];
    return sum;
}
\end{lstlisting}

\begin{lstlisting}[language=bash,caption={Asembler i raport wektoryzacji w GCC (Linux).}]
g++ -std=c++20 -O3 -S -masm=intel dot.cpp -o dot.s

g++ -std=c++20 -O3 -fopt-info-vec-optimized \
    -c dot.cpp -o dot.o
objdump -drC dot.o
\end{lstlisting}

Dla wektoryzacji ważne są: zależności między iteracjami, wyrównanie,
\emph{aliasing} oraz semantyka arytmetyki zmiennoprzecinkowej.
Brak instrukcji SIMD nie musi oznaczać błędu kompilatora; dla
dokładnie zachowanego porządku sumowania kompilator może nie móc
zmienić kolejności operacji. Flaga \code{-ffast-math} pozwala na
istotnie inne przekształcenia numeryczne i~nie powinna być
włączana automatycznie w~programach wymagających ścisłej semantyki.

\subsection{Profile-guided optimization (PGO)}

Kompilator może wykorzystać profile zebrane na reprezentatywnych
danych, aby lepiej rozłożyć kod i~podjąć decyzje optymalizacyjne.
Schemat dla GCC:

\begin{lstlisting}[language=bash,caption={Dwa przebiegi budowania z PGO (GCC).}]
mkdir -p pgo-data
g++ -std=c++20 -O2 -fprofile-generate=pgo-data \
    app.cpp -o app
./app  # Wykonaj reprezentatywny workload.
g++ -std=c++20 -O2 -fprofile-use=pgo-data \
    app.cpp -o app
\end{lstlisting}

PGO nie zastępuje pomiarów: profil z~innego rozkładu danych może
nawet pogorszyć wyniki. Należy też zachować zgodność źródeł, flag
i~miejsc profilu pomiędzy dwoma etapami budowania.

\section{Testowanie kodu wielowątkowego}
\label{sec:cpp18-concurrency-tests}

W~systemach współbieżnych testy wynikowe muszą badać także
\emph{zakończenie pracy}, własność danych i~synchronizację.
Powtórzenie testu wiele razy zwiększa szansę wykrycia błędu, ale
nigdy nie zastępuje rozumowania o~relacjach \emph{happens-before}
(rozdziały~\ref{chap:concurrency} i~\ref{chap:cpp16-advanced-concurrency}).

\begin{lstlisting}[caption={Poprawny test licznika atomowego z jthread (C++20).}]
#include <atomic>
#include <cassert>
#include <thread>
#include <vector>

int main() {
    std::atomic<int> count{0};
    {
        std::vector<std::jthread> workers;
        for (int i = 0; i < 4; ++i) {
            workers.emplace_back([&] {
                for (int k = 0; k < 10000; ++k)
                    count.fetch_add(1, std::memory_order_relaxed);
            });
        }
    } // Destruktory jthread czekaja na wszystkie watki.
    assert(count.load(std::memory_order_relaxed) == 40000);
}
\end{lstlisting}

Do koordynowania momentu rozpoczęcia testu można wykorzystać
\code{std::latch} lub \code{std::barrier} (C++20). W~testach
\emph{timeout} zapobiega nieskończonemu oczekiwaniu w~infrastrukturze
testowej, ale nie może być uznawany za mechanizm anulowania
wątku. Zwracaj uwagę na wyścigi związane z~destrukcją fixture
oraz na kolejność zamykania wątków i~obiektów synchronizacyjnych.

\section{Praktyczny workflow: od błędu do poprawki}

Prześledźmy zalecaną kolejność działania:
\begin{enumerate}
\item \textbf{Odtwórz błąd}: zanotuj wejście, wersję programu, flagi
      kompilacji, środowisko i~opcjonalnie seed.
\item \textbf{Zminimalizuj przypadek}: mały, powtarzalny przykład jest
      łatwiejszy do analizy niż cały system.
\item \textbf{Włącz diagnostykę}: ostrzeżenia, ASan/UBSan, a~dla
      kodu współbieżnego oddzielne uruchomienie TSan.
\item \textbf{Postaw hipotezę}: zidentyfikuj zmienną lub niezmiennik,
      który naruszono; sprawdź stos i~stan programu w~debuggerze.
\item \textbf{Dodaj test regresyjny}: powinien zawodzić przed poprawką
      i~przechodzić po niej.
\item \textbf{Popraw przyczynę}: nie tylko usuń widoczny symptom.
\item \textbf{Zmierz koszt}: jeżeli poprawka ma znaczenie wydajnościowe,
      porównaj benchmarki i~profile przed oraz po zmianie.
\item \textbf{Uruchom pełen zestaw testów}: szczególnie z~różnymi
      konfiguracjami optymalizacji i~kompilatorami.
\end{enumerate}

\section{Automatyzacja w CI i macierz konfiguracji}

W~projektach utrzymywanych długo testy powinny wykonywać się
automatycznie po każdej istotnej zmianie. Praktyczna macierz obejmuje
zwykle konfigurację debugową z~sanitizatorami oraz konfigurację
optymalizowaną z~testami wynikowymi. TSan uruchamiamy w~osobnej
konfiguracji i~na wspieranym środowisku.

\begin{center}
\small
\begin{tabularx}{\textwidth}{@{}p{0.28\textwidth}X@{}}
\toprule
Konfiguracja & Główne zadanie\\
\midrule
Debug, \code{-O0 -g} & Debugowanie źródłowe i~lokalizacja błędów.\\
Debug/san, \code{-O1 -g}, ASan+UBSan & Dostęp do niepoprawnej pamięci i~wybrane UB.\\
Oddzielny build TSan & Wyścigi danych przy konkurencyjnym wykonaniu.\\
Release, \code{-O2} lub \code{-O3} & Testy produkcyjnej konfiguracji i~benchmarki.\\
Clang i GCC & Różne diagnostyki i~różne decyzje optymalizatora.\\
\bottomrule
\end{tabularx}
\end{center}

\nota{Benchmarki są wrażliwe na obciążenie współdzielonych runnerów
CI. W~zwykłym pipeline CI warto przede wszystkim weryfikować
poprawność i~duże regresje; wiarygodne mikrobenchmarki powinny
wykonywać się na kontrolowanym sprzęcie.}

\section{Najczęstsze błędy i wybór narzędzia}

\begin{center}
\small
\begin{tabularx}{\textwidth}{@{}p{0.40\textwidth}X@{}}
\toprule
Problem & Zalecane podejście\\
\midrule
Funkcja daje złe wyniki & Test jednostkowy, debugger, test regresyjny.\\
Awaria tylko dla niektórych danych & Fuzzing, test własnościowy, ASan/UBSan.\\
Podejrzenie wiszącego wskaźnika & ASan, kontrola czasu życia, debugger.\\
Podejrzenie data race & TSan, analiza synchronizacji i~\emph{happens-before}.\\
Program zawiesza się przy wielu wątkach & Stosy wszystkich wątków, graf blokad, timeout testu.\\
Wysokie zużycie pamięci & Massif, profiler alokacji, pomiar RSS.\\
Program jest wolny & \code{perf}/profiler, później benchmark gorącej funkcji.\\
Spadek szybkości po zmianie & Reprodukowalny benchmark, kontrola warunków, profil porównawczy.\\
Wątpliwości co do wektoryzacji & Raport optymalizacji i~asembler, weryfikacja numeryczna.\\
\bottomrule
\end{tabularx}
\end{center}

\section*{Podsumowanie rozdziału}
\addcontentsline{toc}{section}{Podsumowanie rozdziału}

Testy określają, czy otrzymujemy oczekiwane zachowanie; debugger
pozwala prześledzić przyczyny problemu; sanitizatory wykrywają
określone naruszenia reguł wykonywania; profiler i~benchmarki
pomagają podejmować decyzje dotyczące wydajności. Nie istnieje
jedno narzędzie, które zastępuje wszystkie pozostałe.

Praktyczna sekwencja brzmi: \emph{najpierw poprawność, potem
bezpieczeństwo i~reprodukcyjność, następnie pomiary, wreszcie
optymalizacja}. Każda wykryta usterka powinna prowadzić do testu
regresyjnego, a~każda deklarowana poprawa wydajności --- do
porównywalnego pomiaru. Dla systemów wielowątkowych konieczne jest
ponadto świadome rozumienie synchronizacji, czasu życia obiektów
oraz gwarancji modelu pamięci C++.

\section*{Literatura i dokumentacja}
\addcontentsline{toc}{section}{Literatura i dokumentacja}
\begingroup
\sloppy
\begin{itemize}
\item \url{https://google.github.io/googletest/primer.html} ---
      GoogleTest: asercje, fixture i organizacja testów.
\item \url{https://github.com/google/googletest/blob/main/docs/quickstart-cmake.md}
      --- integracja GoogleTest i CMake.
\item \url{https://cmake.org/cmake/help/latest/module/GoogleTest.html}
      --- \code{gtest_discover_tests} i CTest.
\item \url{https://github.com/google/benchmark/blob/main/docs/user_guide.md}
      --- Google Benchmark i stabilne pomiary.
\item \url{https://clang.llvm.org/docs/AddressSanitizer.html} --- ASan.
\item \url{https://clang.llvm.org/docs/UndefinedBehaviorSanitizer.html}
      --- UBSan.
\item \url{https://clang.llvm.org/docs/ThreadSanitizer.html} --- TSan.
\item \url{https://clang.llvm.org/docs/LibFuzzer.html} --- fuzzing.
\item \url{https://sourceware.org/gdb/current/onlinedocs/gdb.html} --- GDB.
\item \url{https://lldb.llvm.org/use/tutorial.html} --- LLDB.
\item \url{https://valgrind.org/docs/manual/manual.html} --- Memcheck i Massif.
\item \url{https://perf.wiki.kernel.org/} --- profiler \code{perf}.
\item Brendan Gregg, \emph{Systems Performance}, wyd. 2,
      Addison-Wesley, 2020.
\end{itemize}
\endgroup


\chapter{Testowanie, debugowanie i profilowanie C++}
\label{chap:cpp18-test-debug-profile}

\section{Trzy różne pytania o program}

Po napisaniu i skompilowaniu programu nie wiemy jeszcze, czy jest on
poprawny, bezpieczny i wydajny. Każdą z tych własności sprawdza się
innymi metodami:
\begin{description}[style=nextline]
\item[Testowanie (\emph{testing})]
Czy program daje oczekiwane wyniki na wybranych danych oraz czy
spełnia wymagane kontrakty? Testy wykrywają błędy, lecz ich pozytywny
wynik nie jest ogólnym dowodem poprawności.
\item[Debugowanie (\emph{debugging})]
Gdzie i dlaczego program odbiega od oczekiwanego zachowania?
Debugger pozwala obserwować stos wywołań, zmienne, wątki i kolejne
instrukcje.
\item[Profilowanie (\emph{profiling})]
Gdzie program zużywa czas procesora, pamięć i inne zasoby?
Profilowanie wskazuje rzeczywiste wąskie gardła; samo zgadywanie
na podstawie długości kodu często prowadzi do błędnych optymalizacji.
\end{description}

Istotne jest także rozróżnienie \emph{testu poprawności} od
\emph{benchmarku}: pierwszy sprawdza warunek logiczny, drugi mierzy
wielkość fizyczną (np. czas, przepustowość, alokacje). Szybki, lecz
błędny algorytm nie jest dobrym rozwiązaniem.

\nota{W~tym rozdziale podstawowym standardem jest C++20. GoogleTest,
Google Benchmark, GDB, LLDB, Valgrind i~narzędzia \code{perf} nie są
częścią standardowej biblioteki C++; wymagają odpowiednich instalacji.
Polecenia zależą od systemu operacyjnego i~kompilatora.}

\section{Strategia testowania i testowalny projekt}
\label{sec:cpp18-strategy}

\subsection{Poziomy testów}

\begin{center}
\small
\begin{tabularx}{\textwidth}{@{}p{0.29\textwidth}X@{}}
\toprule
Rodzaj testu & Co sprawdza i kiedy jest przydatny?\\
\midrule
Jednostkowy (\emph{unit}) & Mała, izolowana funkcja lub klasa; szybka informacja zwrotna.\\
Integracyjny & Współdziałanie komponentów, np. parsera i~bazy danych.\\
Systemowy / end-to-end & Cały scenariusz użytkownika, często z~rzeczywistym wejściem/wyjściem.\\
Regresyjny & Przypadek, który wcześniej ujawniał błąd; nie powinien powrócić.\\
Własnościowy / losowy & Niezmienniki dla wielu zestawów danych, nie tylko pojedyncze wyniki.\\
Wydajnościowy & Czas, przepustowość, opóźnienie, pamięć; wymaga kontrolowanego środowiska.\\
\bottomrule
\end{tabularx}
\end{center}

Dobry test powinien być \emph{powtarzalny}, \emph{niezależny} od innych
testów i~możliwie mały. W~praktyce wydzielamy obliczenia od efektów
ubocznych (sieci, zegara, plików, losowości), dzięki czemu zależności
możemy zastąpić kontrolowanymi implementacjami.

\subsection{Kontrakty, przypadki brzegowe i błędy numeryczne}

Rozważmy funkcję obliczającą średnią. Jej kontrakt powinien ustalać
zachowanie dla pustego zakresu; poniżej zgłaszamy wyjątek.

\begin{lstlisting}[caption={Prosta funkcja z jawnym kontraktem dla pustego zakresu (C++20).}]
#include <numeric>
#include <span>
#include <stdexcept>

double average(std::span<const double> xs) {
    if (xs.empty())
        throw std::invalid_argument("empty range");
    return std::accumulate(xs.begin(), xs.end(), 0.0)
         / static_cast<double>(xs.size());
}
\end{lstlisting}

Testujemy nie tylko wartości typowe, lecz także rozmiary $0$ i~$1$,
wartości ujemne, powtórzenia, skrajne zakresy oraz dane naruszające
warunki wejściowe. Dla liczb zmiennoprzecinkowych zwykła równość
\code{a == b} nie zawsze jest właściwa: wynik może zależeć od porządku
dodawania i~optymalizacji. Stosujemy tolerancję dostosowaną do skali
i~wymagań zadania, a~nie arbitralne \code{1e-6} dla każdego algorytmu.

\section{Testy bez zewnętrznego frameworka}

\subsection{\texttt{assert} i \texttt{static\_assert}}

\code{assert(expr)} z~nagłówka \code{<cassert>} sprawdza warunek
w~czasie działania, lecz może zostać usunięte przy zdefiniowanym
\code{NDEBUG}. \code{static_assert} sprawdza warunek w~czasie
kompilacji. W~szczególności zwykłe \code{assert} nie nadaje się do
walidacji niepoprawnych danych wejściowych w~kodzie produkcyjnym.

\begin{lstlisting}[caption={Asercja kompilacyjna i kontrola zachowania funkcji.}]
#include <cassert>
#include <concepts>
#include <type_traits>

constexpr int square(int x) { return x * x; }
static_assert(square(6) == 36);
static_assert(std::integral<int>);
static_assert(!std::is_copy_constructible_v<void>);

int main() {
    assert(square(7) == 49);
}
\end{lstlisting}

\nota{Przy \code{-DNDEBUG} argument \code{assert} może w~ogóle nie być
wykonywany. Nie wolno umieszczać w~nim niezbędnych efektów ubocznych,
np.~\code{assert(++counter == 1)}. \code{static_assert} również nie
zastępuje testów dla danych znanych dopiero podczas wykonania.}

\subsection{Niezależne testy danych wejściowych}

Poniższą funkcję \code{main} umieść pod definicją \code{average}
z~poprzedniego listingu. Test nie wymaga zewnętrznej biblioteki.
Sprawdza wynik, wyjątek i~przypadek jednoelementowy. Kompiluj go
bez \code{-DNDEBUG}, aby wyrażenia \code{assert} były wykonywane.

\begin{lstlisting}[caption={Test funkcji average bez frameworka; dołącz pod jej definicją (C++20).}]
#include <cassert>
#include <cmath>
#include <stdexcept>
#include <vector>

// Zakladamy definicje average z poprzedniego przykladu.
int main() {
    const std::vector<double> xs{1.0, 2.0, 3.0, 4.0};
    assert(std::abs(average(xs) - 2.5) < 1e-12);
    const std::vector<double> one{7.0};
    assert(average(one) == 7.0);

    bool caught = false;
    const std::vector<double> empty;
    try {
        average(empty);
    } catch (const std::invalid_argument&) {
        caught = true;
    }
    assert(caught);
}
\end{lstlisting}

W~realnym projekcie lepiej użyć frameworka, który rejestruje testy,
raportuje niepowodzenia i~pozwala uruchamiać wybrane przypadki.

\section{GoogleTest: testy jednostkowe w praktyce}
\label{sec:cpp18-gtest}

Biblioteka \emph{GoogleTest} (często \code{gtest}) udostępnia makra
\code{TEST}, \code{TEST_F}, \code{EXPECT_*} i~\code{ASSERT_*}.
Assercje \code{EXPECT_*} zapisują niepowodzenie i~pozwalają wykonać
resztę funkcji testowej, podczas gdy \code{ASSERT_*} przerywają
bieżącą funkcję testową. To nie są mechanizmy obsługi wyjątków C++.

\begin{lstlisting}[caption={GoogleTest: wyniki, wyjatki i przyblizona rownosc.}]
// Plik: average_test.cpp
#include <gtest/gtest.h>
#include <stdexcept>
#include <vector>

#include "average.hpp"
// Funkcja average jest zdefiniowana w pliku average.cpp.
TEST(AverageTest, ComputesMean) {
    const std::vector<double> xs{2.0, 4.0, 6.0};
    EXPECT_NEAR(average(xs), 4.0, 1e-12);
}

TEST(AverageTest, RejectsEmptyInput) {
    const std::vector<double> empty;
    EXPECT_THROW(average(empty), std::invalid_argument);
}

TEST(AverageTest, HandlesSingleton) {
    const std::vector<double> xs{7.0};
    ASSERT_EQ(xs.size(), 1u);
    EXPECT_DOUBLE_EQ(average(xs), 7.0);
}
\end{lstlisting}

\code{EXPECT_NEAR(a,b,tol)} sprawdza bezwzględną różnicę, dlatego
niekiedy lepiej użyć własnego testu tolerancji względnej. Makro
\code{EXPECT_DOUBLE_EQ} porównuje liczby z~tolerancją opartą na ULP;
nie należy utożsamiać go z~matematyczną równością liczb rzeczywistych.

\subsection{Fixture i testy parametryzowane}

\emph{Fixture} przygotowuje stan wspólny dla testów. Każdy test
\code{TEST_F} otrzymuje osobny egzemplarz fixture, dzięki czemu zmiany
jednego testu nie powinny wpływać na inny.

\begin{lstlisting}[caption={Fixture z przygotowanym wektorem.}]
#include <gtest/gtest.h>
#include <algorithm>
#include <vector>

class VectorTest : public ::testing::Test {
protected:
    std::vector<int> values{3, 1, 4, 1, 5};
};

TEST_F(VectorTest, SortsAscending) {
    std::sort(values.begin(), values.end());
    EXPECT_TRUE(std::is_sorted(values.begin(), values.end()));
}

TEST_F(VectorTest, StartsFromOriginalState) {
    EXPECT_EQ(values.front(), 3);
    EXPECT_EQ(values.size(), 5u);
}
\end{lstlisting}

Dla wielu par wejście--wynik przydatne są testy parametryzowane:

\begin{lstlisting}[caption={GoogleTest: test parametryzowany z zestawem wartosci.}]
#include <gtest/gtest.h>
#include <tuple>

constexpr int square(int x) { return x * x; }

class SquareTest
    : public ::testing::TestWithParam<std::tuple<int, int>> {};

TEST_P(SquareTest, MatchesExpected) {
    auto [x, expected] = GetParam();
    EXPECT_EQ(square(x), expected);
}

INSTANTIATE_TEST_SUITE_P(
    Values, SquareTest,
    ::testing::Values(
        std::tuple{0, 0},
        std::tuple{2, 4},
        std::tuple{-3, 9}));
\end{lstlisting}

\subsection{Mocki i testowanie współpracy komponentów}

\emph{GoogleMock} pozwala tworzyć kontrolowane implementacje
interfejsów i~sprawdzać wywołania między obiektami. Używaj go tam,
gdzie zachowanie naprawdę zależy od interakcji; nadmierne mockowanie
utrwala przypadkowe szczegóły implementacji.

\begin{lstlisting}[caption={GoogleMock: sprawdzenie pojedynczego wywolania.}]
#include <gmock/gmock.h>
#include <string_view>

struct Sender {
    virtual ~Sender() = default;
    virtual bool send(std::string_view) = 0;
};

struct MockSender : Sender {
    MOCK_METHOD(bool, send, (std::string_view), (override));
};

void notify(Sender& s) { (void)s.send("ready"); }

TEST(NotificationTest, SendsOneMessage) {
    MockSender sender;
    EXPECT_CALL(sender, send("ready"))
        .Times(1).WillOnce(::testing::Return(true));
    notify(sender);
}
\end{lstlisting}

Własne wyjątki, klasy niewspierające kopiowania i~asynchroniczne
wywołania wymagają odpowiedniego projektu fixture/mocków.
Dla operacji wykonywanych w~innym wątku szczególnie istotne jest
kontrolowane oczekiwanie na koniec pracy przed zniszczeniem mocka.

\subsection{CMake i CTest: mały projekt wieloplikowy}
\label{sec:cpp18-cmake-test}

Poniższa organizacja zakłada, że \code{average} ma deklarację w~pliku
\code{average.hpp}, definicję w~\code{average.cpp}, a~testy w~pliku
\code{average_test.cpp}. W~przykładzie \code{average_test.cpp}
należy dodać \verb|#include "average.hpp"| zamiast deklarować
funkcję ręcznie. Plik nagłówkowy może wyglądać następująco:

\begin{lstlisting}[caption={Plik average.hpp: deklaracja API testowanej funkcji.}]
#pragma once
#include <span>
double average(std::span<const double> xs);
\end{lstlisting}

Plik \code{average.cpp} zawiera kod pierwszego przykładu
z~tej części rozdziału (z~dołączonym \code{average.hpp});
w~pliku testów umieszczamy testy z~sekcji~\ref{sec:cpp18-gtest}.

\begin{lstlisting}[language=bash,caption={Struktura projektu z GoogleTest.}]
project/
  CMakeLists.txt
  average.hpp
  average.cpp
  average_test.cpp
\end{lstlisting}

\begin{lstlisting}[language=bash,caption={CMakeLists.txt: biblioteka, testy i automatyczna rejestracja w CTest.}]
cmake_minimum_required(VERSION 3.20)
project(cpp18_tests LANGUAGES CXX)

add_library(core average.cpp)
target_compile_features(core PUBLIC cxx_std_20)
target_include_directories(core PUBLIC ${CMAKE_CURRENT_SOURCE_DIR})

# GoogleTest musi byc zainstalowany lub dostarczony przez menedzer pakietow.
find_package(GTest CONFIG REQUIRED)

add_executable(core_tests average_test.cpp)
target_link_libraries(core_tests PRIVATE core GTest::gtest_main)

enable_testing()
include(GoogleTest)
gtest_discover_tests(core_tests)
\end{lstlisting}

\begin{lstlisting}[language=bash,caption={Budowanie i uruchamianie testow przez CMake/CTest.}]
cmake -S . -B build -DCMAKE_BUILD_TYPE=Debug
cmake --build build -j
ctest --test-dir build --output-on-failure
./build/core_tests --gtest_filter='AverageTest.*'
\end{lstlisting}

\nota{W~systemach z~generatorem wielokonfiguracyjnym (np.~Visual
Studio) wariant budowania wybiera się przez \code{--config Debug}.
Przykład wymaga zainstalowanego GoogleTest; alternatywą jest
\code{FetchContent} z~\emph{przypiętą} wersją zależności, a~nie
pobieranie nieokreślonej wersji podczas każdego budowania.}

\section{Testy własnościowe, generatywne i fuzzing}
\label{sec:cpp18-property}

Sprawdzanie setek wartości jest bardziej użyteczne, jeśli testuje
\emph{własność}, a~nie jedynie listę ręcznie wpisanych wyników.
Na przykład po posortowaniu ciągu musi on być niemalejący i~mieć
tę samą multizbiorowość wartości co przed sortowaniem.

\begin{lstlisting}[caption={Powtarzalny test wlasnosciowy przy uzyciu generatora losowego.}]
#include <algorithm>
#include <cassert>
#include <random>
#include <vector>

int main() {
    std::mt19937 gen(123456u);  // Ustalony seed.
    std::uniform_int_distribution<int> count(0, 300);
    std::uniform_int_distribution<int> value(-1000, 1000);

    for (int trial = 0; trial < 1000; ++trial) {
        std::vector<int> x(count(gen));
        for (int& v : x) v = value(gen);
        auto original = x;
        std::sort(x.begin(), x.end());
        assert(std::is_sorted(x.begin(), x.end()));
        std::sort(original.begin(), original.end());
        assert(x == original);
    }
}
\end{lstlisting}

W~testach losowych przechowuj \emph{seed} i~minimalny przykład,
który ujawnił błąd. Warto również stosować \emph{shrinking}:
automatyczne zmniejszanie wejścia do możliwie prostego kontrprzykładu.

\emph{Fuzzing} automatycznie generuje wejścia, kierując się m.in.
pokryciem kodu. W~Clangu narzędzie \emph{libFuzzer} może wywoływać
poniższą funkcję z~wieloma binarnymi danymi wejściowymi.

\begin{lstlisting}[caption={Minimalny punkt wejsciowy dla libFuzzer (Clang).}]
#include <cstddef>
#include <cstdint>
#include <vector>

extern "C" int LLVMFuzzerTestOneInput(
    const std::uint8_t* data, std::size_t size) {
    // Podmien na testowany parser lub dekoder.
    std::vector<std::uint8_t> bytes(data, data + size);
    if (bytes.size() > 2 && bytes[0] == 0x42) {
        // Przypadek zalezy od danych; nie wolno zakladac ich poprawnosci.
        (void)bytes[1];
    }
    return 0;
}
\end{lstlisting}

\begin{lstlisting}[language=bash,caption={Budowanie harnessu fuzzingu w Clangu.}]
clang++ -std=c++20 -O1 -g \
  -fsanitize=fuzzer,address,undefined fuzz.cpp -o fuzz
./fuzz -max_total_time=30
\end{lstlisting}

Po zgłoszeniu awarii zapisujemy minimalny przypadek jako test
regresyjny. Fuzzer nie dowodzi braku błędów, a~jego skuteczność zależy
od jakości punktów wejściowych i~stanu pokrycia.

\section{Debugowanie: od komunikatu do przyczyny}
\label{sec:cpp18-debug}

\subsection{Informacje debugowe i rodzaje kompilacji}

Użyteczne opcje GCC/Clang to \code{-g} (symbole debugowe),
\code{-O0} (łatwe śledzenie krokowe) oraz \code{-Og} (optymalizacje
starające się zachować dobrą obserwowalność w~debuggerze). Program
kompilowany z~\code{-O2} może nie mieć osobnej zmiennej ani instrukcji
odpowiadającej każdej linii kodu źródłowego.

\begin{lstlisting}[language=bash,caption={Budowanie programu do debugowania.}]
g++ -std=c++20 -g -O0 -Wall -Wextra -Wpedantic \
    app.cpp -o app_debug

gdb ./app_debug
\end{lstlisting}

Nie wszystkie błędy są odtwarzalne przy \code{-O0}; niektóre pojawiają
się dopiero przy \code{-O2} lub \code{-O3}. Do ich analizy warto
pozostawić \code{-g} także w~wersji zoptymalizowanej i~zachować
niestripowane symbole debugowe.

\subsection{GDB: breakpoint, stos i obserwacja zmiennych}

\begin{lstlisting}[caption={Program do cwiczenia debugowania.}]
#include <stdexcept>
#include <vector>

int element_at(const std::vector<int>& v, std::size_t index) {
    return v.at(index);  // Moze zglosic std::out_of_range.
}

int main() {
    const std::vector<int> v{3, 5, 8};
    return element_at(v, 5);
}
\end{lstlisting}

\begin{lstlisting}[language=bash,caption={Typowa sesja w GDB.}]
(gdb) break element_at        # Punkt zatrzymania.
(gdb) run                     # Uruchom program.
(gdb) info args               # Argumenty funkcji.
(gdb) print index             # Wartosc zmiennej.
(gdb) backtrace               # Stos wywolan.
(gdb) next                    # Nastepna linia bez wchodzenia.
(gdb) step                    # Wejdz do wywolywanej funkcji.
(gdb) catch throw             # Zatrzymaj przy zglaszaniu wyjatku.
(gdb) continue                # Kontynuuj.
\end{lstlisting}

\code{break}, \code{watch}, \code{next}, \code{step},
\code{finish}, \code{frame}, \code{print} i~\code{bt} to podstawowe
narzędzia. \emph{Watchpoint} śledzi zmianę pamięci, ale możliwość
sprzętowego śledzenia zależy od architektury i~debuggera.

\subsection{Wątki i deadlock w debuggerze}

W~aplikacji wielowątkowej sama analiza wątku \code{main} nie
wystarczy. W~GDB:

\begin{lstlisting}[language=bash,caption={Diagnostyka wielowatkowosci w GDB.}]
(gdb) info threads
(gdb) thread 3
(gdb) bt
(gdb) thread apply all bt
(gdb) set print pretty on
\end{lstlisting}

Jeśli wątki zatrzymały się w~\code{mutex::lock()} lub
\code{condition_variable::wait()}, trzeba odtworzyć graf
zależności: który wątek posiada dany mutex i~na co czeka.
Program zakleszczony nie musi zawierać wyścigu danych, więc
ThreadSanitizer nie zastąpi analizy zakleszczeń.

\subsection{LLDB: odpowiednik na macOS i innych platformach}

LLDB oferuje analogiczne operacje:

\begin{lstlisting}[language=bash,caption={Podstawowe polecenia LLDB.}]
(lldb) breakpoint set --name element_at
(lldb) run
(lldb) frame variable
(lldb) thread backtrace
(lldb) thread list
(lldb) thread backtrace all
(lldb) next
(lldb) step
(lldb) continue
\end{lstlisting}

W~VS Code można podłączyć GDB/LLDB poprzez odpowiednie rozszerzenie,
a~następnie korzystać z~breakpointów, podglądu zmiennych i~stosów.
Dokładny adapter oraz konfiguracja \code{launch.json} zależą od
systemu, rozszerzenia i~narzędzi zainstalowanych lokalnie.

\section{Sanitizatory i narzędzia wykrywania błędów}
\label{sec:cpp18-sanitizers}

Sanitizatory dodają do programu kod kontrolujący określone klasy
błędów w~czasie wykonywania. W~odróżnieniu od debuggera mogą
wykryć nieprawidłową operację \emph{w~chwili jej wykonania}, zanim
zauważymy błędny rezultat.

\begin{center}
\small
\begin{tabularx}{\textwidth}{@{}p{0.28\textwidth}X@{}}
\toprule
Narzędzie & Główne zastosowanie\\
\midrule
AddressSanitizer (ASan) & Wyjście poza granice, use-after-free, część przypadków wycieków.\\
UndefinedBehaviorSanitizer (UBSan) & Np. przepełnienie \code{signed int}, błędne przesunięcia, null dereference, alignment.\\
ThreadSanitizer (TSan) & Wiele wyścigów danych i~część problemów synchronizacji.\\
MemorySanitizer (MSan) & Odczyty niezainicjalizowanej pamięci; wymaga specjalnych warunków instrumentacji.\\
Valgrind Memcheck & Wykrywanie wybranych błędów pamięci i~wycieków (na wspieranych systemach).\\
\bottomrule
\end{tabularx}
\end{center}

\subsection{AddressSanitizer i UndefinedBehaviorSanitizer}

Poniższy program jest \emph{celowo niepoprawny}; dostęp do elementu
po zwolnieniu pamięci jest niezdefiniowanym zachowaniem. Nie należy
interpretować jego wyniku bez instrumentacji.

\begin{lstlisting}[caption={Celowo bledny program: use-after-free.}]
int main() {
    int* p = new int(42);
    delete p;
    return *p;  // UB: odczyt z nieistniejacego obiektu.
}
\end{lstlisting}

\begin{lstlisting}[language=bash,caption={Instrumentacja ASan i UBSan; GCC/Clang.}]
g++ -std=c++20 -O1 -g -fno-omit-frame-pointer \
    -fsanitize=address,undefined -fno-sanitize-recover=all \
    bad.cpp -o bad_san
./bad_san
\end{lstlisting}

ASan i~UBSan można często połączyć. ASan/TSan/MSan nie są natomiast
zwykle łączone ze sobą w~jednym procesie: używamy osobnych
konfiguracji i~uruchomień. Raporty sanitizatorów zależą od wersji
kompilatora, wsparcia runtime i~platformy.

\subsection{ThreadSanitizer i wyścig danych}

Przykładowy wyścig występuje przy równoczesnym modyfikowaniu
zwykłego licznika. To również \emph{celowo błędny} kod:

\begin{lstlisting}[caption={Celowo bledne wspoldzielenie licznika.}]
#include <thread>

int main() {
    int x = 0;
    std::thread t1([&] { ++x; });
    std::thread t2([&] { ++x; });
    t1.join();
    t2.join();
}
\end{lstlisting}

\begin{lstlisting}[language=bash,caption={Oddzielna konfiguracja z ThreadSanitizer.}]
g++ -std=c++20 -O1 -g -pthread \
    -fsanitize=thread race.cpp -o race_tsan
./race_tsan
\end{lstlisting}

\nota{Brak zgłoszenia w~danym uruchomieniu \emph{nie} dowodzi braku
wyścigów danych. TSan ma ograniczenia platformowe, może zgłaszać
wyniki fałszywie dodatnie dla niestandardowej synchronizacji i~nie
gwarantuje wykrywania deadlocków.}

\subsection{Valgrind i testy wycieków}

Na wspieranych systemach Linux przydatna jest kontrola programów
bez kompilacyjnej instrumentacji:

\begin{lstlisting}[language=bash,caption={Przykladowe badanie pamieci przez Valgrind Memcheck.}]
valgrind --tool=memcheck --leak-check=full \
         --show-leak-kinds=all ./app_debug
\end{lstlisting}

Różne narzędzia mają inne modele wykrywania i~inne koszty
wykonania. Wybór zależy od konkretnej klasy błędów; wyniki dobrze
weryfikować na małym, odtwarzalnym przypadku testowym.

\section{Analiza statyczna, ostrzeżenia i pokrycie kodu}
\label{sec:cpp18-static}

Kompilator potrafi zgłosić wiele usterek jeszcze przed wykonaniem
programu. Warto utrzymywać ostrzeżenia w~CI i~unikać wyciszania ich
bez zrozumienia przyczyny.

\begin{lstlisting}[language=bash,caption={Przykladowe ostrzezenia GCC/Clang.}]
g++ -std=c++20 -Wall -Wextra -Wpedantic -Wconversion \
    -Wshadow -Werror -c core.cpp
\end{lstlisting}

Opcje zależą od kompilatora; \code{-Werror} bywa pożyteczne przy
stabilnej bazie kodu, lecz aktualizacja kompilatora może wprowadzić
nowe diagnostyki. \emph{clang-tidy} wykonuje zestaw analiz
statycznych, np. dotyczących nowoczesnych idiomów C++, błędów
potencjalnych i~stylu.

\begin{lstlisting}[language=bash,caption={Analiza projektu ze zbiorem komend kompilacji.}]
cmake -S . -B build -DCMAKE_EXPORT_COMPILE_COMMANDS=ON
clang-tidy -p build src/core.cpp
\end{lstlisting}

\subsection{Pokrycie testami: co liczby naprawdę oznaczają?}

\emph{Line coverage} mierzy wykonane linie, \emph{branch coverage}
wykonane gałęzie, a~\emph{function coverage} wywołane funkcje.
Wysoki procent nie oznacza sprawdzenia wszystkich kombinacji wejść
ani dowodu poprawności. Na przykład test może wykonać gałąź, nie
sprawdzając rezultatu.

\begin{lstlisting}[language=bash,caption={Pokrycie testow z GCC i gcov (przykladowy projekt).}]
g++ -std=c++20 --coverage -O0 -g core.cpp tests.cpp -o tests
./tests
gcov core.cpp
\end{lstlisting}

Dla Clanga dostępny jest osobny zestaw instrumentacji LLVM:
\begin{lstlisting}[language=bash,caption={Clang: instrumentacja i raport pokrycia.}]
clang++ -std=c++20 -O0 -g -fprofile-instr-generate \
    -fcoverage-mapping core.cpp tests.cpp -o tests
LLVM_PROFILE_FILE=tests.profraw ./tests
llvm-profdata merge -sparse tests.profraw -o tests.profdata
llvm-cov report ./tests -instr-profile=tests.profdata
\end{lstlisting}

W~szczególności należy testować \emph{warunki brzegowe}, a~nie
maksymalizować jedynie procent pokrycia.

\section{Benchmarking: jak mierzyć szybkość bez samooszukiwania}
\label{sec:cpp18-benchmark}

\subsection{Pułapki mikrobenchmarków}

Nowoczesny kompilator usuwa obliczenia o~nieobserwowalnych efektach,
wykonuje \emph{constant folding} i~wektoryzuje pętle. Dlatego mierzenie
pustej pętli lub stałego obliczenia może nie mieć związku z~realną
wydajnością. Wynik zależy ponadto od rozgrzania cache, alokacji,
obciążenia systemu, zegarów CPU, temperatury i~danych testowych.

\begin{lstlisting}[caption={Pomiar algorytmu z rzeczywistym wynikiem (bez biblioteki benchmark).}]
#include <chrono>
#include <iostream>
#include <numeric>
#include <vector>

int main() {
    std::vector<int> data(1'000'000, 1);
    using clock = std::chrono::steady_clock;
    long long checksum = 0;

    for (int repetition = 0; repetition < 10; ++repetition) {
        auto start = clock::now();
        const auto result = std::accumulate(
            data.begin(), data.end(), 0LL);
        auto finish = clock::now();
        checksum += result;  // Wynik jest wykorzystany po pomiarze.
        auto us = std::chrono::duration_cast<
            std::chrono::microseconds>(finish - start).count();
        std::cout << "Trial " << repetition << ": " << us << " us\n";
    }
    std::cout << "checksum=" << checksum << '\n';
}
\end{lstlisting}

To ilustracja, nie wzorcowy mikrobenchmark: optymalizator może
ułatwić sobie obliczenia na regularnych danych; dodatkowo
\code{std::cout} poza odcinkiem pomiarowym zmienia warunki
między próbami. Użyteczne benchmarki wymagają wielu powtórzeń,
przygotowanych danych i~analizy rozkładu czasów.

\subsection{Google Benchmark}

\emph{Google Benchmark} automatyzuje wielokrotne uruchamianie i~raportuje
miary czasowe. W~krytycznym fragmencie wykorzystujemy
\code{benchmark::DoNotOptimize}, aby wynik obliczenia nie został
trywialnie usunięty; nie wyłącza to jednak wszystkich optymalizacji
na badanym wyrażeniu.

\begin{lstlisting}[caption={Samodzielny benchmark sumowania wielu elementow.}]
// Plik: benchmark_sum.cpp
#include <benchmark/benchmark.h>
#include <numeric>
#include <vector>

static void BM_Accumulate(benchmark::State& state) {
    const std::vector<int> data(state.range(0), 1);
    for (auto _ : state) {
        const auto sum = std::accumulate(
            data.begin(), data.end(), 0LL);
        benchmark::DoNotOptimize(sum);
    }
    state.SetItemsProcessed(state.iterations() * state.range(0));
}

BENCHMARK(BM_Accumulate)->Arg(1024)->Arg(1 << 20);
BENCHMARK_MAIN();
\end{lstlisting}

\begin{lstlisting}[language={},caption={CMakeLists.txt dla przykładu Google Benchmark.}]
cmake_minimum_required(VERSION 3.20)
project(BenchmarkSum LANGUAGES CXX)
find_package(benchmark CONFIG REQUIRED)
add_executable(bench benchmark_sum.cpp)
target_compile_features(bench PRIVATE cxx_std_20)
target_link_libraries(bench PRIVATE benchmark::benchmark)
\end{lstlisting}

\begin{lstlisting}[language=bash,caption={Budowanie i powtarzanie pomiaru.}]
cmake -S . -B build -DCMAKE_BUILD_TYPE=Release
cmake --build build
./build/bench --benchmark_repetitions=10
./build/bench --benchmark_out=results.json \
              --benchmark_out_format=json
\end{lstlisting}

W~porównaniach równoległych trzeba ponadto ustalić liczbę wątków,
rozmiary bloków, warm-up, powtarzalność danych oraz ewentualne
\emph{false sharing}. Sam najkrótszy pojedynczy czas nie opisuje
stabilnej wydajności.

\section{Profilowanie CPU i pamięci}
\label{sec:cpp18-profiling}

\subsection{Profilowanie próbkowe: Linux \texttt{perf}}

Zamiast mierzyć wyłącznie całkowity czas, profiler pomaga wykrywać
\emph{hot spots}: funkcje zajmujące znaczącą część wykonania.
Przykładowe polecenia systemu Linux:

\begin{lstlisting}[language=bash,caption={Profilowanie CPU przez perf (Linux).}]
g++ -std=c++20 -O2 -g -fno-omit-frame-pointer \
    app.cpp -o app
perf stat ./app
perf record -g ./app
perf report
\end{lstlisting}

\code{perf stat} podaje liczniki zdarzeń udostępniane przez jądro
i~sprzęt, np. cykle i~instrukcje. \code{perf record} zbiera próbki
umożliwiające oglądanie gorących funkcji i~stosu. Dostępne zdarzenia
i~uprawnienia zależą od komputera i~ustawień systemu.

\subsection{Flame graph}

\emph{Flame graph} to agregat stosów wywołań: szerokość prostokąta
odpowiada udziałowi próbek, a~wysokość zagnieżdżeniu wywołań.
Nie jest to dokładna oś czasu. Przykładowy przepływ z~zewnętrznymi
narzędziami FlameGraph:

\begin{lstlisting}[language=bash,caption={Tworzenie flame graph z probek perf (Linux).}]
perf record -F 99 -g -- ./app
perf script | stackcollapse-perf.pl | flamegraph.pl > flame.svg
\end{lstlisting}

Ostatnie dwa skrypty pochodzą z~osobnego projektu FlameGraph i~muszą
być zainstalowane. W~projektach o~dużej liczbie wątków pomagają też
profile opóźnień, I/O i~blokad, a~nie tylko wykorzystania CPU.

\subsection{Alokacje i zużycie pamięci}

Wydajność pamięci to nie tylko liczba bajtów. Ważne są rozkład
alokacji, fragmentacja, lokalność cache i~czas życia obiektów.
Na wspieranych systemach można użyć m.in.:

\begin{lstlisting}[language=bash,caption={Przykladowe pomiary alokacji i pamieci.}]
valgrind --tool=massif ./app
ms_print massif.out.*
/usr/bin/time -v ./app
\end{lstlisting}

\code{massif} służy do profilowania sterty (zgodnie z~możliwościami
narzędzia); \code{/usr/bin/time -v} raportuje miary procesu, np.
\emph{maximum resident set size}, ale nie mówi, które alokacje były
odpowiedzialne. Instrumentacja pamięci również wpływa na czas pracy.

\section{Asembler i związek kodu źródłowego z wydajnością}
\label{sec:cpp18-asm}

Dwa podobne zapisy C++ mogą dać różny kod maszynowy, szczególnie po
włączeniu \code{-O2} lub \code{-O3}. Przy analizie trzeba jednak
uważać, aby nie wyciągać wniosków na podstawie sztucznych przykładów,
w których kompilator może wyeliminować całą pracę.

\begin{lstlisting}[caption={Niewielka funkcja do analizy wygenerowanego kodu.}]
#include <cstddef>

double dot(const double* a, const double* b, std::size_t n) {
    double sum = 0.0;
    for (std::size_t i = 0; i < n; ++i)
        sum += a[i] * b[i];
    return sum;
}
\end{lstlisting}

\begin{lstlisting}[language=bash,caption={Asembler i raport wektoryzacji w GCC (Linux).}]
g++ -std=c++20 -O3 -S -masm=intel dot.cpp -o dot.s

g++ -std=c++20 -O3 -fopt-info-vec-optimized \
    -c dot.cpp -o dot.o
objdump -drC dot.o
\end{lstlisting}

Dla wektoryzacji ważne są: zależności między iteracjami, wyrównanie,
\emph{aliasing} oraz semantyka arytmetyki zmiennoprzecinkowej.
Brak instrukcji SIMD nie musi oznaczać błędu kompilatora; dla
dokładnie zachowanego porządku sumowania kompilator może nie móc
zmienić kolejności operacji. Flaga \code{-ffast-math} pozwala na
istotnie inne przekształcenia numeryczne i~nie powinna być
włączana automatycznie w~programach wymagających ścisłej semantyki.

\subsection{Profile-guided optimization (PGO)}

Kompilator może wykorzystać profile zebrane na reprezentatywnych
danych, aby lepiej rozłożyć kod i~podjąć decyzje optymalizacyjne.
Schemat dla GCC:

\begin{lstlisting}[language=bash,caption={Dwa przebiegi budowania z PGO (GCC).}]
mkdir -p pgo-data
g++ -std=c++20 -O2 -fprofile-generate=pgo-data \
    app.cpp -o app
./app  # Wykonaj reprezentatywny workload.
g++ -std=c++20 -O2 -fprofile-use=pgo-data \
    app.cpp -o app
\end{lstlisting}

PGO nie zastępuje pomiarów: profil z~innego rozkładu danych może
nawet pogorszyć wyniki. Należy też zachować zgodność źródeł, flag
i~miejsc profilu pomiędzy dwoma etapami budowania.

\section{Testowanie kodu wielowątkowego}
\label{sec:cpp18-concurrency-tests}

W~systemach współbieżnych testy wynikowe muszą badać także
\emph{zakończenie pracy}, własność danych i~synchronizację.
Powtórzenie testu wiele razy zwiększa szansę wykrycia błędu, ale
nigdy nie zastępuje rozumowania o~relacjach \emph{happens-before}
(rozdziały~\ref{chap:concurrency} i~\ref{chap:cpp16-advanced-concurrency}).

\begin{lstlisting}[caption={Poprawny test licznika atomowego z jthread (C++20).}]
#include <atomic>
#include <cassert>
#include <thread>
#include <vector>

int main() {
    std::atomic<int> count{0};
    {
        std::vector<std::jthread> workers;
        for (int i = 0; i < 4; ++i) {
            workers.emplace_back([&] {
                for (int k = 0; k < 10000; ++k)
                    count.fetch_add(1, std::memory_order_relaxed);
            });
        }
    } // Destruktory jthread czekaja na wszystkie watki.
    assert(count.load(std::memory_order_relaxed) == 40000);
}
\end{lstlisting}

Do koordynowania momentu rozpoczęcia testu można wykorzystać
\code{std::latch} lub \code{std::barrier} (C++20). W~testach
\emph{timeout} zapobiega nieskończonemu oczekiwaniu w~infrastrukturze
testowej, ale nie może być uznawany za mechanizm anulowania
wątku. Zwracaj uwagę na wyścigi związane z~destrukcją fixture
oraz na kolejność zamykania wątków i~obiektów synchronizacyjnych.

\section{Praktyczny workflow: od błędu do poprawki}

Prześledźmy zalecaną kolejność działania:
\begin{enumerate}
\item \textbf{Odtwórz błąd}: zanotuj wejście, wersję programu, flagi
      kompilacji, środowisko i~opcjonalnie seed.
\item \textbf{Zminimalizuj przypadek}: mały, powtarzalny przykład jest
      łatwiejszy do analizy niż cały system.
\item \textbf{Włącz diagnostykę}: ostrzeżenia, ASan/UBSan, a~dla
      kodu współbieżnego oddzielne uruchomienie TSan.
\item \textbf{Postaw hipotezę}: zidentyfikuj zmienną lub niezmiennik,
      który naruszono; sprawdź stos i~stan programu w~debuggerze.
\item \textbf{Dodaj test regresyjny}: powinien zawodzić przed poprawką
      i~przechodzić po niej.
\item \textbf{Popraw przyczynę}: nie tylko usuń widoczny symptom.
\item \textbf{Zmierz koszt}: jeżeli poprawka ma znaczenie wydajnościowe,
      porównaj benchmarki i~profile przed oraz po zmianie.
\item \textbf{Uruchom pełen zestaw testów}: szczególnie z~różnymi
      konfiguracjami optymalizacji i~kompilatorami.
\end{enumerate}

\section{Automatyzacja w CI i macierz konfiguracji}

W~projektach utrzymywanych długo testy powinny wykonywać się
automatycznie po każdej istotnej zmianie. Praktyczna macierz obejmuje
zwykle konfigurację debugową z~sanitizatorami oraz konfigurację
optymalizowaną z~testami wynikowymi. TSan uruchamiamy w~osobnej
konfiguracji i~na wspieranym środowisku.

\begin{center}
\small
\begin{tabularx}{\textwidth}{@{}p{0.28\textwidth}X@{}}
\toprule
Konfiguracja & Główne zadanie\\
\midrule
Debug, \code{-O0 -g} & Debugowanie źródłowe i~lokalizacja błędów.\\
Debug/san, \code{-O1 -g}, ASan+UBSan & Dostęp do niepoprawnej pamięci i~wybrane UB.\\
Oddzielny build TSan & Wyścigi danych przy konkurencyjnym wykonaniu.\\
Release, \code{-O2} lub \code{-O3} & Testy produkcyjnej konfiguracji i~benchmarki.\\
Clang i GCC & Różne diagnostyki i~różne decyzje optymalizatora.\\
\bottomrule
\end{tabularx}
\end{center}

\nota{Benchmarki są wrażliwe na obciążenie współdzielonych runnerów
CI. W~zwykłym pipeline CI warto przede wszystkim weryfikować
poprawność i~duże regresje; wiarygodne mikrobenchmarki powinny
wykonywać się na kontrolowanym sprzęcie.}

\section{Najczęstsze błędy i wybór narzędzia}

\begin{center}
\small
\begin{tabularx}{\textwidth}{@{}p{0.40\textwidth}X@{}}
\toprule
Problem & Zalecane podejście\\
\midrule
Funkcja daje złe wyniki & Test jednostkowy, debugger, test regresyjny.\\
Awaria tylko dla niektórych danych & Fuzzing, test własnościowy, ASan/UBSan.\\
Podejrzenie wiszącego wskaźnika & ASan, kontrola czasu życia, debugger.\\
Podejrzenie data race & TSan, analiza synchronizacji i~\emph{happens-before}.\\
Program zawiesza się przy wielu wątkach & Stosy wszystkich wątków, graf blokad, timeout testu.\\
Wysokie zużycie pamięci & Massif, profiler alokacji, pomiar RSS.\\
Program jest wolny & \code{perf}/profiler, później benchmark gorącej funkcji.\\
Spadek szybkości po zmianie & Reprodukowalny benchmark, kontrola warunków, profil porównawczy.\\
Wątpliwości co do wektoryzacji & Raport optymalizacji i~asembler, weryfikacja numeryczna.\\
\bottomrule
\end{tabularx}
\end{center}

\section*{Podsumowanie rozdziału}
\addcontentsline{toc}{section}{Podsumowanie rozdziału}

Testy określają, czy otrzymujemy oczekiwane zachowanie; debugger
pozwala prześledzić przyczyny problemu; sanitizatory wykrywają
określone naruszenia reguł wykonywania; profiler i~benchmarki
pomagają podejmować decyzje dotyczące wydajności. Nie istnieje
jedno narzędzie, które zastępuje wszystkie pozostałe.

Praktyczna sekwencja brzmi: \emph{najpierw poprawność, potem
bezpieczeństwo i~reprodukcyjność, następnie pomiary, wreszcie
optymalizacja}. Każda wykryta usterka powinna prowadzić do testu
regresyjnego, a~każda deklarowana poprawa wydajności --- do
porównywalnego pomiaru. Dla systemów wielowątkowych konieczne jest
ponadto świadome rozumienie synchronizacji, czasu życia obiektów
oraz gwarancji modelu pamięci C++.

\section*{Literatura i dokumentacja}
\addcontentsline{toc}{section}{Literatura i dokumentacja}
\begingroup
\sloppy
\begin{itemize}
\item \url{https://google.github.io/googletest/primer.html} ---
      GoogleTest: asercje, fixture i organizacja testów.
\item \url{https://github.com/google/googletest/blob/main/docs/quickstart-cmake.md}
      --- integracja GoogleTest i CMake.
\item \url{https://cmake.org/cmake/help/latest/module/GoogleTest.html}
      --- \code{gtest_discover_tests} i CTest.
\item \url{https://github.com/google/benchmark/blob/main/docs/user_guide.md}
      --- Google Benchmark i stabilne pomiary.
\item \url{https://clang.llvm.org/docs/AddressSanitizer.html} --- ASan.
\item \url{https://clang.llvm.org/docs/UndefinedBehaviorSanitizer.html}
      --- UBSan.
\item \url{https://clang.llvm.org/docs/ThreadSanitizer.html} --- TSan.
\item \url{https://clang.llvm.org/docs/LibFuzzer.html} --- fuzzing.
\item \url{https://sourceware.org/gdb/current/onlinedocs/gdb.html} --- GDB.
\item \url{https://lldb.llvm.org/use/tutorial.html} --- LLDB.
\item \url{https://valgrind.org/docs/manual/manual.html} --- Memcheck i Massif.
\item \url{https://perf.wiki.kernel.org/} --- profiler \code{perf}.
\item Brendan Gregg, \emph{Systems Performance}, wyd. 2,
      Addison-Wesley, 2020.
\end{itemize}
\endgroup
